\documentclass[11pt]{amsart}
\usepackage{amsmath,amssymb,amsthm,booktabs,hyperref,url,array}
\usepackage[margin=1.1in]{geometry}
\newtheorem{theorem}{Theorem}[section]
\newtheorem{proposition}[theorem]{Proposition}
\newtheorem{conjecture}[theorem]{Conjecture}
\newtheorem{lemma}[theorem]{Lemma}
\newtheorem{corollary}[theorem]{Corollary}
\theoremstyle{definition}
\newtheorem{definition}[theorem]{Definition}
\newtheorem{observation}[theorem]{Observation}
\newtheorem{remark}[theorem]{Remark}
\newtheorem{example}[theorem]{Example}
\newcommand{\Q}{\mathbb{Q}}
\newcommand{\Z}{\mathbb{Z}}
\newcommand{\F}{\mathbb{F}}
\newcommand{\T}{\mathbb{T}}
\newcommand{\degphi}{\deg\varphi_E}
\newcommand{\ord}{\operatorname{ord}}
\newcommand{\Phig}{\Phi^{\mathrm{geom}}}
\newcommand{\Gal}{\operatorname{Gal}}
\newcommand{\ad}{\operatorname{ad}^0}
\newcommand{\Sym}{\operatorname{Sym}^2}
\newcommand{\nv}{\mathrm{nv}}
\newcommand{\Frob}{\operatorname{Frob}}

\title[Local factors of the adjoint and the modular degree]{Local factors of the adjoint $L$-function and the $\ell$-adic valuations of the modular degree}
\author{Dharmaj Soni}
\email{dharmaj.soni@gmail.com}
\urladdr{https://orcid.org/0009-0007-6397-3351}
\date{\today}

\begin{document}

\begin{abstract}
Let $E/\Q$ be the $X_0(N)$-optimal curve of its isogeny class, $\varphi_E\colon X_0(N)\to E$ its modular parametrisation and $f$ its newform. For an odd prime $\ell$ at which $E[\ell]$ is irreducible we study $v_\ell(\degphi)$ through the adjoint $L$-function $L(\ad f,s)=L(\Sym f,s+1)$. The congruence number of $f$, which is $\degphi$ up to primes whose square divides $N$, is governed by the \emph{naive} adjoint $L$-value, whose Euler product omits the local factors at the additive primes. We conjecture that $v_\ell(\degphi)$ is at least the sum, over the bad primes $q\ne\ell$, of the $\ell$-adic valuation of the omitted factor at $s=1$, plus the Tamagawa exponent of $\ad$ at the potentially multiplicative primes (the $\ell$-part of the valuation of the Tate parameter), the terms of the Kodaira types $\mathrm{III}$ and $\mathrm{III}^*$ set aside. For a tame prime $q\ge5$ the omitted factor contributes $v_\ell(q-\chi_{-3}(q))$ for the types $\mathrm{II},\mathrm{II}^*,\mathrm{IV},\mathrm{IV}^*$, $v_\ell(q^2-1)$ for $\mathrm{I}_n^*$ and $v_\ell(q-1)+v_\ell((q+1)^2-a_q(E')^2)$ for $\mathrm{I}_0^*$, with $E'$ the twist of good reduction. The inequality holds with no exception, with the Euler factors computed by PARI at every bad prime, for every optimal curve of conductor at most $500000$ and for $893$ curves of conductor up to $4\cdot10^6$ outside the tables, at $\ell=3,5,7,11,13$, with equality in $57\%$ of the cases with a positive bound at $\ell=3$. We prove the terms of the types $\mathrm{I}_0^*$ and $\mathrm{I}_n^*$ for every odd $\ell\nmid 2N$ with $E[\ell]$ irreducible, from a formula of Diamond, Flach and Guo for the naive congruence ideal and an elementary twisting argument; the Tamagawa terms for $\ell\ge5$ from a theorem of Kim and Ota; and the qualitative statements for the other types at $\ell=3$ by level lowering. The rest is reduced to a statement about the adjoint Selmer group, and the type $\mathrm{III}$ terms are exactly those for which that statement fails in the data, by one, at $\ell=3$ only. The inequality by local component groups and its twist-symmetrised refinement are corollaries, and we describe the failure in the Eisenstein case.
\end{abstract}

\maketitle

\section{Introduction}

Let $E/\Q$ be an elliptic curve of conductor $N$ and let $f=f_E\in S_2(\Gamma_0(N))$ be the newform attached to it. Among the curves in the isogeny class of $E$ there is a distinguished one, the $X_0(N)$-optimal curve, for which the modular parametrisation $\varphi_E\colon X_0(N)\to E$ has minimal degree; throughout, $E$ denotes this curve and $\degphi$ its \emph{modular degree}. The modular degree is computable \cite{Cremona97,Watkins02}, and its size is tied to deep questions (the conjectured bound $\degphi\ll_\varepsilon N^{2+\varepsilon}$ is equivalent to a form of the $abc$ conjecture, see \cite[\S2]{ARS12} and the references there), but its prime factorisation is not understood. The purpose of this note is to describe, and partly prove, a local-to-global divisibility: the $\ell$-adic valuation of $\degphi$ is bounded below by local invariants of $E$ at the bad primes, and the right invariants are the local factors of the adjoint $L$-function.

The classical route to the degree is the symmetric square. Up to explicit factors, $\degphi$ is the special value $L(\Sym f,2)$ divided by a period \cite{Flach92,Watkins02}, and by a theorem of Agashe, Ribet and Stein \cite[Theorem 2.2]{ARS12} it divides the \emph{congruence number} $r_E$, the largest integer $r$ such that $f\equiv g\pmod r$ coefficientwise for some integral cusp form $g$ of level $N$ orthogonal to $f$, with $v_p(\degphi)=v_p(r_E)$ for every prime $p$ with $p^2\nmid N$. The congruence number is the congruence ideal of $f$ in the full Hecke algebra of level $N$ (Lemma~\ref{lem:crit}), and the theory of such ideals, from Hida \cite{Hida81} to Diamond, Flach and Guo \cite{DFG04}, expresses them through the adjoint Selmer group and the adjoint $L$-value. The $L$-function that appears there is the \emph{naive} one: the Euler product of $\sum a_n(f)^2n^{-s}$, which has the true local factor at the primes of good or multiplicative reduction and the trivial factor at the additive primes. The point of this note is that the primes at which the naive and the true adjoint $L$-functions differ are exactly the primes at which the modular degree acquires extra local divisibility, and that the amount is the valuation of the omitted factor.

\subsection{The local terms}
Let $V=V_\ell E$, $\rho=\rho_{E,\ell}$ and $A_f=\ad V$, the trace-zero adjoint. We follow the conventions of \cite[\S1.7.2]{DFG04}: for a prime $p$ let $\delta_p=\dim V^{I_p}$, so $\delta_p=2$, $1$, $0$ for good, multiplicative, additive reduction; $L_p(A_f,s)=\det(1-\Frob_p\,p^{-s}\mid A_f^{I_p})^{-1}$ is the true local factor, and the naive factor is $L_p^{\nv}(A_f,s)=L_p(A_f,s)$ if $\delta_p>0$ and $1$ if $\delta_p=0$. In the normalisation of the symmetric square, $L_p(\Sym f,s)=L_p(A_f,s-1)$, and PARI's \texttt{lfunsympow(E,2)} returns these factors at every prime, including $2$ and $3$ \cite{PARI}.

\begin{definition}\label{def:terms}
For a bad prime $q\ne\ell$ put
\[
e_\ell(E,q)=v_\ell\Bigl(\frac{L_q(A_f,1)^{-1}}{L_q^{\nv}(A_f,1)^{-1}}\Bigr)=v_\ell\Bigl(\frac{P_q(q^{-2})}{1-a_q(E)^2q^{-2}}\Bigr),\quad
t_\ell(E,q)=\begin{cases}v_\ell(-v_q(j_E))&\text{if }v_q(j_E)<0,\\0&\text{otherwise,}\end{cases}
\]
where $P_q(X)$ is the polynomial with $L_q(\Sym f,s)=P_q(q^{-s})^{-1}$. So $e_\ell$ is the valuation of the Euler factor that the naive $L$-function omits; it is $0$ at a multiplicative prime, where the naive factor is the true one. The term $t_\ell$ is the Tamagawa exponent of $\ad V$ at a potentially multiplicative prime, the $\ell$-part of the valuation of the Tate parameter; it is $v_\ell(n)$ for a Kodaira symbol $\mathrm{I}_n$ at any $q$ and $\mathrm{I}_n^*$ at odd $q$, while a symbol $\mathrm{I}_n^*$ at $q=2$ can have potentially good reduction, and then $t_\ell=0$. Both terms are invariant under quadratic twist.
\end{definition}

For a tame prime $q\ge5$ the omitted factor is read off the Kodaira type (Lemma~\ref{lem:euler}). With $\chi_{-3}$, $\chi_{-4}$ the quadratic characters of $\Q(\sqrt{-3})$ and $\Q(i)$, $E'$ the quadratic twist of $E$ by $\Q(\sqrt{q^*})$, $q^*=(-1)^{(q-1)/2}q$, and $\alpha,\beta$ the roots of $X^2-a_q(E')X+q$ (for type $\mathrm{I}_0^*$, where $E'$ has good reduction at $q$):
\begin{center}\footnotesize\setlength{\tabcolsep}{4pt}
\begin{tabular}{lll}
\toprule
type at $q\ge5$ & $L_q(A_f,s)^{-1}$ & $e_\ell(E,q)$\\
\midrule
$\mathrm{I}_n$ & $1-q^{-1-s}$ & $0$\\
$\mathrm{I}_n^*$ & $1-q^{-1-s}$ & $v_\ell(q^2-1)$\\
$\mathrm{II},\mathrm{II}^*,\mathrm{IV},\mathrm{IV}^*$ & $1-\chi_{-3}(q)q^{-s}$ & $v_\ell(q-\chi_{-3}(q))$\\
$\mathrm{III},\mathrm{III}^*$ & $1-\chi_{-4}(q)q^{-s}$ & $v_\ell(q-\chi_{-4}(q))$\\
$\mathrm{I}_0^*$ & $(1-\tfrac{\alpha}{\beta}q^{-s})(1-q^{-s})(1-\tfrac{\beta}{\alpha}q^{-s})$ & $v_\ell(q-1)+v_\ell\bigl((q+1)^2-a_q(E')^2\bigr)$\\
\bottomrule
\end{tabular}
\end{center}

\subsection{The rule}
\begin{conjecture}[Euler-factor rule]\label{conj:euler}
Let $E$ be $X_0(N)$-optimal and without complex multiplication, and let $\ell$ be an odd prime such that $E[\ell]$ is irreducible as a $\Gal(\overline\Q/\Q)$-module. Then
\begin{equation}\label{eq:rule}
v_\ell(\degphi)\ \ge\ \sum_{\substack{q\mid N,\ q\ne\ell\\ \text{type}\ne\mathrm{III},\mathrm{III}^*}}e_\ell(E,q)\ +\ \sum_{q\mid N}t_\ell(E,q)\ +\ [\ell=3]\,w_3(E),
\end{equation}
where $w_3(E)$ is the weight of the prime $3$ itself when $3\mid N$ and the reduction at $3$ is potentially good: $1$ for the types $\mathrm{II},\mathrm{IV},\mathrm{III}^*$, $2$ for $\mathrm{II}^*,\mathrm{IV}^*$ and $0$ for $\mathrm{III},\mathrm{I}_0^*$.
\end{conjecture}

Three comments on the statement. The weights $w_3$ at $q=\ell=3$ are empirical (Section~\ref{sec:wild}); at $q=\ell\ge5$ no term is needed. The terms at the wild primes $2$ and $3$ are the ones of Definition~\ref{def:terms} with PARI's Euler factor, and they reproduce every minimum observed in the tables, including the irregular behaviour of the prime $2$ noted in an earlier version of this note (Section~\ref{sec:wild}). The Euler terms of the types $\mathrm{III}$ and $\mathrm{III}^*$ are attained in isolation like all the others (Table~\ref{tab:iso}), and they can be added to \eqref{eq:rule} except, at $\ell=3$ only, when such a prime meets a Tamagawa term at a prime $p\equiv-1\pmod3$: that configuration fails by exactly one in about $2\%$ of the cases, and it is the test case for the Selmer-group statement of Section~\ref{sec:selmer}.

\subsection{Local component groups}
The rule contains the two statements of the earlier version of this note. For a bad prime $q$ let $\Phig_q=\Phi_q(\overline\F_q)$ be the group of components of the special fibre of the N\'eron model of $E$ at $q$ over the algebraic closure. Its order is read off the Kodaira type: $n$ for type $\mathrm{I}_n$; $1$ for $\mathrm{II},\mathrm{II}^*$; $2$ for $\mathrm{III},\mathrm{III}^*$; $3$ for $\mathrm{IV},\mathrm{IV}^*$; $4$ for $\mathrm{I}_0^*$ and $\mathrm{I}_n^*$ \cite[IV.9]{SilvermanATAEC}. The Tamagawa number $c_q=\#\Phi_q(\F_q)$ counts only the components defined over $\F_q$; the statements below are about $\Phig_q$, not $c_q$. For an odd prime $q$ let $E^{(q)}$ be the quadratic twist of $E$ by $\Q(\sqrt{q^*})$; it has the same reduction type as $E$ at every prime other than $q$, while at $q$ the Kodaira types are exchanged (Lemma~\ref{lem:twist}):
\[
\mathrm{I}_n\leftrightarrow\mathrm{I}_n^*,\qquad \mathrm{II}\leftrightarrow\mathrm{IV}^*,\qquad \mathrm{IV}\leftrightarrow\mathrm{II}^*,\qquad \mathrm{III}\leftrightarrow\mathrm{III}^*,\qquad \mathrm{I}_0\leftrightarrow\mathrm{I}_0^*
\]
(the first and the last for every odd $q$, the three others for $q\ge5$). Define the \emph{twist-symmetrised} weight
\[
w_\ell(E,q)=\begin{cases}
v_\ell(n) & \text{if the type at $q$ is $\mathrm{I}_n$, or $\mathrm{I}_n^*$ with $q$ odd},\\
1 & \text{if $\ell=3$ and the type is $\mathrm{IV}$ or $\mathrm{IV}^*$, or $\mathrm{II}$ or $\mathrm{II}^*$ with $q$ odd},\\
0 & \text{otherwise},
\end{cases}
\]
so that
\[
w_\ell(E,q)=\max\bigl(v_\ell(\#\Phig_q(E)),\ v_\ell(\#\Phig_q(E^{(q)}))\bigr)\ \text{for }q\ge5,\qquad w_\ell(E,2)=v_\ell(\#\Phig_2(E)).
\]
The two inequalities are
\begin{equation}\label{eq:main}
v_\ell(\degphi)\ \ge\ \sum_{q\mid N} v_\ell\bigl(\#\Phig_q\bigr),
\end{equation}
\begin{equation}\label{eq:twist}
v_\ell(\degphi)\ \ge\ \sum_{q\mid N} w_\ell(E,q),
\end{equation}
of which the second implies the first. Both were conjectured in the earlier version under the hypotheses of Conjecture~\ref{conj:euler}, and both hold with no exception on the same ranges (Section~\ref{sec:forms}).

\begin{proposition}\label{prop:forms}
Assume Conjecture~\ref{conj:euler} for $(E,\ell)$. Then \eqref{eq:main} and \eqref{eq:twist} hold, unless $\ell=3$ and $E$ has type $\mathrm{IV}$ or $\mathrm{IV}^*$ at $2$, in which case they hold as soon as $e_3(E,2)\ge1$; that inequality holds for every such curve of the tables.
\end{proposition}
\begin{proof}
Compare the two sides term by term; $w_\ell(E,q)\ge v_\ell(\#\Phig_q)$, so it suffices to treat \eqref{eq:twist}. At a prime of type $\mathrm{I}_n$ both sides have $v_\ell(n)=t_\ell(E,q)$. At a prime of type $\mathrm{I}_n^*$ with $q$ odd, $t_\ell(E,q)=v_\ell(n)=w_\ell(E,q)$; with $q=2$, $w_\ell(E,2)=v_\ell(4)=0$. At $q\ge5$ of type $\mathrm{II},\mathrm{II}^*,\mathrm{IV},\mathrm{IV}^*$ and $\ell=3$, $e_3(E,q)=v_3(q-\chi_{-3}(q))\ge1=w_3(E,q)$ since $q\equiv\chi_{-3}(q)\pmod3$; for $\ell\ge5$ both weights are $0$. At $q=3$ of these types and $\ell=3$, $w_3(E)\ge1=w_3(E,3)$ by definition, and for $\ell\ge5$ the weight is $0$. At $q=2$ of type $\mathrm{II},\mathrm{II}^*$, $w_\ell(E,2)=0$; of type $\mathrm{IV},\mathrm{IV}^*$ and $\ell=3$, $w_3(E,2)=1$, which is the excepted case. The types $\mathrm{III},\mathrm{III}^*,\mathrm{I}_0^*$ have $w_\ell=0$ for odd $\ell$. All other terms of \eqref{eq:rule} are nonnegative.
\end{proof}

\subsection{What is proved}
The multiplicative part of \eqref{eq:main} has a long history. Ribet and Takahashi \cite{RibetTakahashi97,Takahashi01} compared $\degphi$ with the degree of a parametrisation by a Shimura curve and found the product of the $n_q$ over the primes of the quaternion discriminant, up to primes where $E[\ell]$ is reducible; Pasten \cite{Pasten} controls that error term. Pollack and Weston \cite{PollackWeston11} conjectured, and proved for weight $2$ and square-free level when $\bar\rho$ is ramified at two primes at least, and Kim and Ota \cite{KimOta} proved in general, that moving a prime $q$ into the ``new'' part of the level changes the congruence ideal of $f_E$ in the full space $S_2(\Gamma_0(N))$ (for an elliptic curve, the ideal generated by $r_E$) by exactly the \emph{Tamagawa exponent} $t_q$, the largest $t$ with $\rho_{E,\ell}\bmod\ell^t$ unramified at $q$; for multiplicative $q$ one has $t_q=v_\ell(n_q)=t_\ell(E,q)$ whether or not the reduction is split (Lemma~\ref{lem:tamagawa-exp}). With the theorem of Agashe, Ribet and Stein this gives:

\begin{theorem}[Kim--Ota, Agashe--Ribet--Stein]\label{thm:known}
Let $\ell\ge5$ be a prime not dividing $N$ such that $\bar\rho_{E,\ell}$ is absolutely irreducible when restricted to $\Gal(\overline\Q/\Q(\sqrt{\ell^*}))$, $\ell^*=(-1)^{(\ell-1)/2}\ell$. Then
\[
v_\ell(\degphi)\ \ge \sum_{\substack{q\,\|\,N\\ q\not\equiv\pm1\ (\mathrm{mod}\ \ell)}} v_\ell(n_q).
\]
\end{theorem}

The data of Section~\ref{sec:data} say that none of the three restrictions in Theorem~\ref{thm:known} is needed for the inequality: it holds at $\ell=3$, at primes $\ell$ dividing $N$, and with the primes $q\equiv\pm1\pmod\ell$ included, in every one of more than two million optimal curves. The qualitative half of the additive terms at $\ell=3$ is a theorem:

\begin{theorem}\label{thm:qual}
Let $E/\Q$ have conductor $N$ and let $\ell$ be an odd prime not dividing $N$ such that $\bar\rho_{E,\ell}$ is irreducible and, if $\ell=3$, absolutely irreducible when restricted to $\Gal(\overline\Q/\Q(\sqrt{-3}))$. Let $q\neq\ell$ be a prime at which $E$ has additive reduction of one of the following kinds:
\begin{enumerate}
\item[(a)] type $\mathrm{IV}$ or $\mathrm{IV}^*$, with $q\ge5$ and $\ell=3$;
\item[(b)] type $\mathrm{II}$ or $\mathrm{II}^*$, with $q\ge5$ and $\ell=3$;
\item[(c)] type $\mathrm{I}_n^*$ with $\ell\mid n$, $q$ odd.
\end{enumerate}
Then $\ell$ divides the modular degree of the $X_0(N)$-optimal curve in the isogeny class of $E$.
\end{theorem}

In each case the prime $q$ has a positive term in \eqref{eq:rule}, and the theorem says that it contributes at least one to $v_\ell(\degphi)$. In case (a) the mechanism is that the ramification of $\bar\rho_{E,3}$ at $q$ is unipotent, so that the level can be lowered from $q^2$ to $q$ (Lemma~\ref{lem:IV}). In cases (b) and (c) $\bar\rho_{E,\ell}$ itself has conductor exponent $2$ at $q$ and nothing can be lowered; but its twist $\bar\rho_{E,\ell}\otimes\chi_q$, the representation of $E^{(q)}$, has conductor exponent $1$ or $0$ at $q$. Lowering the level of the twist and twisting back produces an eigenform of level dividing $N$, different from $f_E$ and congruent to it modulo $\ell$ (Section~\ref{sec:proof}). The quantitative statement is a theorem for the two types whose twist lowers the level:

\begin{theorem}\label{thm:quant}
Let $\ell$ be an odd prime with $\ell\nmid 2N$ and $E[\ell]$ irreducible, and let $S$ be the set of primes $q\ge5$ at which $E$ has type $\mathrm{I}_0^*$ or $\mathrm{I}_n^*$. Then
\[
v_\ell(\degphi)\ \ge\ \sum_{q\in S}e_\ell(E,q).
\]
If moreover $\ell\ge5$ and $\bar\rho_{E,\ell}$ is absolutely irreducible on $\Gal(\overline\Q/\Q(\sqrt{\ell^*}))$, let $M'$ be the set of primes $q$ at which $E$ has type $\mathrm{I}_n$, or $\mathrm{I}_n^*$ with $q$ odd, with $\ell\mid n$ and $q\not\equiv\pm1\pmod\ell$. Then
\[
v_\ell(\degphi)\ \ge\ \sum_{q\in S}e_\ell(E,q)+\sum_{q\in M'}t_\ell(E,q).
\]
\end{theorem}

The first inequality (Theorem~\ref{thm:euler}) is where the adjoint $L$-function enters: the twist $E_S$ of $E$ by the characters of the primes of $S$ has lower conductor, $f_E$ is the $S$-depleted form of $f_{E_S}$, and the formula of Diamond, Flach and Guo for the naive $\Sigma$-finite congruence ideal \cite[Proposition 1.4]{DFG04} says that depleting at a prime multiplies the congruence ideal by the Euler factor at that prime. An elementary twisting argument (Theorem~\ref{thm:twist}) carries the congruences back to level $N$. The only inputs are that formula, which needs $\bar\rho$ irreducible and $\ell\nmid2N$, an algebraic inequality from \cite[Lemma 4.17]{DDT}, and multiplicity one at the lower level; no Taylor--Wiles hypothesis enters, and $\ell=3$ is allowed. The second inequality adds the theorem of Kim and Ota, applied to the twist. For the potentially good types $\mathrm{II}$, $\mathrm{IV}$, $\mathrm{III}$ and their duals the twist does not lower the level and the omitted factor sits inside the naive $L$-value; Section~\ref{sec:selmer} reformulates the rule for them as an inequality between the adjoint Selmer group and the local terms, which is where the question now stands.

The last statement concerns the Eisenstein case, in which \eqref{eq:main} can fail.

\begin{conjecture}\label{conj:eis}
Let $E$ be optimal, not CM, with a rational $\ell$-isogeny, $\ell$ odd, and let $D$ be the deficit $\sum_q v_\ell(\#\Phig_q)-v_\ell(\degphi)$ of \eqref{eq:main}. Then
\[
D\ \le\ v_\ell\bigl(\#E(\Q)_{\mathrm{tors}}\bigr)+\bigl[\,E[3]\simeq\Z/3\oplus\mu_3\,\bigr],
\]
where $[\,\cdot\,]$ is $1$ if the condition holds and $0$ otherwise, and also $D\le$ the number of curves in the isogeny class of $E$ with a rational point of order $\ell$. In particular $D\le0$ when $E$ has no rational $\ell$-torsion point, $D\le1$ for $\ell\ge5$, and $D\le2$ for $\ell=3$, since a point of order $9$ and the image $\Z/3\oplus\mu_3$ do not occur together in the tables (Section~\ref{sec:eis}).
\end{conjecture}

Section~\ref{sec:local} collects the local facts, Section~\ref{sec:proof} proves Theorem~\ref{thm:qual}, Section~\ref{sec:quant} proves Theorem~\ref{thm:quant}, Section~\ref{sec:selmer} contains the Selmer-group reformulation, Section~\ref{sec:data} describes the computations and Section~\ref{sec:two} discusses the prime $2$, where the analogous statements overlap with the literature on Watkins's conjecture.

\section{Local invariants}\label{sec:local}

Let $q$ be a prime of bad reduction. We use Tate's algorithm \cite[IV.9]{SilvermanATAEC} for the Kodaira type and $c_q$, and the following facts.

\begin{lemma}\label{lem:tamagawa-exp}
Let $q$ be a prime of multiplicative reduction, $n_q=v_q(\Delta_E)$, and $\ell\ne q$ a prime. Then $\rho_{E,\ell}\bmod\ell^t$ is unramified at $q$ if and only if $\ell^t\mid n_q$. This holds for split and for nonsplit reduction.
\end{lemma}

\begin{proof}
Over the unramified quadratic extension of $\Q_q$ (which does not change the inertia group) $E$ is a Tate curve with parameter $u$, $v_q(u)=n_q$. The $\ell^t$-torsion is generated by $\mu_{\ell^t}$ and an $\ell^t$-th root of $u$, so inertia acts trivially on $E[\ell^t]$ exactly when $u$ is an $\ell^t$-th power in $\Q_q^{\mathrm{ur}}$ up to units, that is, when $\ell^t\mid n_q$.
\end{proof}

\begin{lemma}\label{lem:IV}
Let $q\ge5$ be a prime at which $E$ has reduction of type $\mathrm{IV}$ or $\mathrm{IV}^*$. Then the image of the inertia group $I_q$ in $\operatorname{Aut}(T_3E)$ is cyclic of order $3$, and its image in $\operatorname{Aut}(E[3])$ is a nontrivial unipotent subgroup. Consequently $E[3]^{I_q}$ is a line, $\bar\rho_{E,3}$ is tamely ramified at $q$, and the Serre conductor of $\bar\rho_{E,3}$ has exponent $1$ at $q$.
\end{lemma}

\begin{proof}
For $q\ge5$ the reduction is potentially good and the semistability defect, the order of the image of $I_q$ in $\operatorname{Aut}(T_\ell E)$ for any $\ell\ne q$, equals $12/\gcd(v_q(\Delta_E),12)$ \cite{Kraus90}; for $v_q(\Delta_E)\in\{4,8\}$ this is $3$. The image is tame since $q\ne3$, hence cyclic, generated by an element $\sigma$ of order $3$ with determinant $1$, so with eigenvalues the two primitive cube roots of unity. Hence $\sigma^2+\sigma+1$ annihilates $T=T_3E$, and $T$ is a module over $\Z_3[\sigma]/(\sigma^2+\sigma+1)\simeq\Z_3[\zeta_3]$, a discrete valuation ring with uniformiser $\pi=1-\zeta_3$ and residue field $\F_3$. Being torsion-free of $\Z_3$-rank $2$, $T$ is free of rank one over $\Z_3[\zeta_3]$, and $E[3]=T/3T\simeq\Z_3[\zeta_3]/(\pi^2)$ with $\sigma$ acting as multiplication by $\zeta_3=1-\pi$. Since $\pi\ne0$ in $\Z_3[\zeta_3]/(\pi^2)$, this is a nontrivial unipotent automorphism of order $3$; its fixed space is the line $\pi\Z_3[\zeta_3]/(\pi^2)$. The Serre conductor exponent at $q$ is $\dim E[3]-\dim E[3]^{I_q}=1$, there being no wild part.
\end{proof}

\begin{remark}\label{rem:types}
For the other additive types at $q\ge5$ the image of inertia mod $3$ is not unipotent: it has order $2$ or $4$ (types $\mathrm{I}_0^*$, $\mathrm{III}$, $\mathrm{III}^*$, defects $2$ and $4$), or it is minus a unipotent (types $\mathrm{II}$, $\mathrm{II}^*$, defect $6$), or it is a ramified quadratic character times a unipotent (type $\mathrm{I}_n^*$). In the last two cases $E[3]^{I_q}=0$ and the Serre conductor of $\bar\rho_{E,3}$ has exponent $2$ at $q$, the same as the conductor of $E$; these are the types whose twist by $\chi_q$ lowers the conductor exponent.
\end{remark}

\begin{lemma}\label{lem:twist}
Let $q$ be an odd prime, $\chi=\chi_q$ the quadratic character of conductor $q$ and $E^{(q)}$ the twist of $E$ by $\chi$, that is, by $\Q(\sqrt{q^*})$.
\begin{enumerate}
\item[(i)] At every prime $p\ne q$ the curves $E$ and $E^{(q)}$ have the same Kodaira type, the same geometric component group and the same conductor exponent.
\item[(ii)] If $E$ has type $\mathrm{I}_n$ at $q$ ($n\ge0$) then $E^{(q)}$ has type $\mathrm{I}_n^*$ at $q$, and conversely, with the same $n$. If $q\ge5$ and $E$ has type $\mathrm{II}$, $\mathrm{III}$, $\mathrm{IV}$ at $q$ then $E^{(q)}$ has type $\mathrm{IV}^*$, $\mathrm{III}^*$, $\mathrm{II}^*$ respectively, and conversely. The conductor exponent of $E^{(q)}$ at $q$ is $2$ when its reduction is additive and $1$ when it is of type $\mathrm{I}_n$, $n\ge1$.
\item[(iii)] (Atkin--Li \cite{AtkinLi78}) Let $g$ be a newform of weight $2$ on $\Gamma_0(M)$ with $q^2\nmid M$. Then the twist $g\otimes\chi$, with coefficients $a_n(g\otimes\chi)=\chi(n)a_n(g)$, is a newform of weight $2$ on $\Gamma_0(M')$, $M'=\operatorname{lcm}(M,q^2)$; that is, $M'=Mq^2$ if $q\nmid M$ and $M'=Mq$ if $q\,\|\,M$.
\end{enumerate}
\end{lemma}

\begin{proof}
(i) The character $\chi$ is unramified at $p$, so $E$ and $E^{(q)}$ become isomorphic over an unramified extension of $\Q_p$, and the Kodaira type, the geometric component group and the conductor exponent are invariant under unramified base change. (ii) If $(a_1,a_2,a_3,a_4,a_6)$ is a minimal model of $E$ at $q$ with invariants $c_4,c_6,\Delta$, the twist has a model with invariants $q^{2}c_4$, $q^{3}c_6$, $q^{6}\Delta$ and the same $j$-invariant. For potentially multiplicative reduction, $v_q(j)<0$, the type is $\mathrm{I}_n$ or $\mathrm{I}_n^*$ with $n=-v_q(j)$ according as the reduction is multiplicative or additive, and a ramified quadratic twist exchanges the two. For $n=0$: a model $y^2=x^3+a_2x^2+a_4x+a_6$ with good reduction at $q$ has twist $y^2=x^3+q^*a_2x^2+q^{*2}a_4x+q^{*3}a_6$, for which the cubic $T^3+a_2T^2+a_4T+a_6$ of Tate's algorithm has distinct roots modulo $q$, the criterion for type $\mathrm{I}_0^*$; conversely a curve of type $\mathrm{I}_0^*$ at an odd $q$ has a model of the latter shape, and its twist has good reduction. For $q\ge5$ and potentially good reduction the type is determined by $v_q(\Delta)$ of a minimal model, which lies in $\{0,2,3,4,6,8,9,10\}$ for the types $\mathrm{I}_0,\mathrm{II},\mathrm{III},\mathrm{IV},\mathrm{I}_0^*,\mathrm{IV}^*,\mathrm{III}^*,\mathrm{II}^*$; the twisted model has $v_q(\Delta)$ increased by $6$, and is minimal exactly when the result is less than $12$, otherwise $12$ is subtracted. The conductor exponents are the standard ones for $q$ odd (tame). (iii) is \cite[Theorem 3.1]{AtkinLi78}; in the case $q\,\|\,M$, the local component of $g$ at $q$ is an unramified twist of the Steinberg representation, whose twist by a ramified quadratic character has conductor exponent $2$.
\end{proof}

\begin{lemma}[the omitted factors at tame primes]\label{lem:euler}
Let $q\ge5$ be a bad prime, $q\ne\ell$. Then $L_q(A_f,s)^{-1}$ and $e_\ell(E,q)$ are as in the table of Section~1.
\end{lemma}

\begin{proof}
The inertia group acts on $V$ through a finite cyclic quotient when the reduction is potentially good, and through a quadratic character times a unipotent when it is potentially multiplicative; $L_q(A_f,s)^{-1}=\prod(1-\lambda q^{-s})$ over the eigenvalues $\lambda$ of Frobenius on $(\ad V)^{I_q}$, in the normalisation in which a good prime has the eigenvalues $\alpha/\beta,1,\beta/\alpha$ on $\ad V$.

\emph{Type $\mathrm{I}_n$.} The local component of $f$ is an unramified twist of the Steinberg representation; $(\ad V)^{I_q}$ is a line with eigenvalue $q^{-1}$, so $L_q(A_f,s)^{-1}=1-q^{-1-s}$, which is also the naive factor since $a_q^2=1$: $e_\ell=0$. \emph{Type $\mathrm{I}_n^*$.} $V=V'\otimes\chi_q$ with $V'$ Steinberg at $q$ (Lemma~\ref{lem:twist}), $\ad V=\ad V'$ and $a_q=0$, so the naive factor is $1$ and $e_\ell=v_\ell(1-q^{-2})=v_\ell(q^2-1)$. \emph{Types $\mathrm{II},\mathrm{II}^*,\mathrm{IV},\mathrm{IV}^*$.} The image of $I_q$ in $\operatorname{Aut}V$ is cyclic of order $3$ or $6$ \cite{Kraus90}, generated by $\sigma$ or $-\sigma$ with $\sigma$ of order $3$ and eigenvalues $\zeta_3^{\pm1}$. On $\ad V$ the sign disappears and the invariants form the line spanned by $\sqrt{-3}\in\Z_\ell[\zeta_3]\simeq\operatorname{End}_{I_q}(V)$; Frobenius centralises $\sigma$ if $q\equiv1\pmod3$ and inverts it if $q\equiv2$, so it acts on the line by $\chi_{-3}(q)$, and $e_\ell=v_\ell(1-\chi_{-3}(q)q^{-1})=v_\ell(q-\chi_{-3}(q))$. \emph{Types $\mathrm{III},\mathrm{III}^*$.} The image has order $4$, eigenvalues $\pm i$; the invariant line of $\ad V$ is spanned by the diagonal element, on which Frobenius acts by $\chi_{-4}(q)$, and $e_\ell=v_\ell(q-\chi_{-4}(q))$. \emph{Type $\mathrm{I}_0^*$.} $V=V'\otimes\chi_q$ with $V'$ unramified with Frobenius eigenvalues $\alpha,\beta$; $\ad V$ is unramified with eigenvalues $\alpha/\beta,1,\beta/\alpha$, and
\[
L_q(A_f,1)^{-1}=\Bigl(1-\frac{\alpha}{\beta q}\Bigr)\Bigl(1-\frac1q\Bigr)\Bigl(1-\frac{\beta}{\alpha q}\Bigr)=\frac{q-1}{q}\cdot\frac{(q+1)^2-a_q(E')^2}{q^2},
\]
using $\alpha\beta=q$ and $\alpha+\beta=a_q(E')$. In each case the values were confirmed against PARI's \texttt{lfunsympow} on every curve of the tables.
\end{proof}

\begin{remark}\label{rem:wild}
At $q=2,3$ the Kodaira type does not determine $L_q(A_f,s)$, but PARI does, and Definition~\ref{def:terms} applies verbatim: for instance a type $\mathrm{IV}$ at $2$ has $P_2(X)=1+2X$ and $e_3=v_3(1+\tfrac12)=1$, a type $\mathrm{II}$ at $2$ has $P_2=1$ and $e_3=0$, and a type $\mathrm{I}_n^*$ at $2$ has potentially good reduction in about half of the cases, with a cubic $P_2$ and $t_\ell=0$. At $q=\ell$ the Euler factor is not the relevant quantity, and the weights $w_3$ of Conjecture~\ref{conj:euler} are empirical.
\end{remark}

\section{Proof of Theorem~\ref{thm:qual}}\label{sec:proof}

Throughout this section $\ell$ is odd and $\ell\nmid N$. Let $f=f_E\in S_2(\Gamma_0(N))$ be the newform of the isogeny class of $E$, let $E_0$ be the optimal curve in the class, and let $r=r_{E_0}$ be its congruence number: the largest integer $r$ such that $f\equiv g\pmod r$ coefficientwise for some $g\in S_2(\Gamma_0(N))$ with integer coefficients orthogonal to $f$ \cite[\S2]{ARS12}. By \cite[Theorem 2.2]{ARS12}, $\deg\varphi_{E_0}$ divides $r$ and $v_p(\deg\varphi_{E_0})=v_p(r)$ for every prime $p$ with $p^2\nmid N$; as $\ell\nmid N$ it suffices to prove $\ell\mid r$. Let $\T$ be the Hecke algebra of level $N$, the subring of $\operatorname{End}S_2(\Gamma_0(N))$ generated over $\Z$ by all $T_n$, $n\ge1$, including the operators $U_p$ for $p\mid N$, and write $\bar\rho=\bar\rho_{E,\ell}$, irreducible by hypothesis. The two lemmas below reduce the theorem to the production of a congruent newform of lower level.

\begin{lemma}[congruence criterion]\label{lem:crit}
Let $\phi_f\colon\T\to\Z$, $T_n\mapsto a_n(f)$, let $I=\ker\phi_f$ and let $J$ be the annihilator in $\T$ of the orthogonal complement $f^\perp$ of $f$. Then $\T/(I+J)$ is cyclic of order $r$; equivalently, $r$ generates the congruence ideal $\phi_f(J)=\phi_f(\operatorname{Ann}_\T I)$. Consequently, if there is a normalised eigenform $g\in S_2(\Gamma_0(N))\otimes\overline\Q$ for all of $\T$, orthogonal to $f$, with coefficients in the ring of integers $\mathcal{O}$ of a number field, and a prime $\lambda\mid\ell$ of $\mathcal{O}$ such that $a_n(g)\equiv a_n(f)\pmod\lambda$ for every $n\ge1$, then $\ell\mid r$.
\end{lemma}

\begin{proof}
Writing $\T\otimes\Q=\Q\times\T'$ with $\T'$ acting faithfully on $f^\perp$ (the $f$-line is $\T$-stable, and so is $f^\perp$, because $f$ is an eigenvector of the adjoints of the $T_n$ as well), one has $I=\T\cap(0\times\T')$ and $J=\T\cap(\Q\times0)$, so $I\cap J=0$ and $\T/(I+J)\simeq\Z/\phi_f(J)$ is cyclic, of order $r'$ say. The pairing $S_2(\Gamma_0(N);\Z)\times\T\to\Z$, $(g,t)\mapsto a_1(tg)$, is perfect \cite[Theorem 2.2]{Ribet83}, \cite[Lemma 2.1]{AU96} (see \cite[\S3]{ARS12}), and under it the forms killed by $I$ are $\Z f$ and the forms killed by $J$ are $L=f^\perp\cap S_2(\Gamma_0(N);\Z)$. Dualising the inclusion $\T\subset\T/I\times\T/J$, whose cokernel is $\T/(I+J)$, shows that $\Z f\oplus L$ has index $r'$ in $S_2(\Gamma_0(N);\Z)$. The quotient $S_2(\Gamma_0(N);\Z)/(\Z f\oplus L)$ embeds in $\Q f/\Z f$ by projection onto $\Q f$ along $f^\perp$, so it is cyclic, generated by the class of a form $h$ with projection $f/r'$; then $r'h=f+g_0$ with $g_0\in L$, so $f\equiv-g_0\pmod{r'}$ and $r'\le r$; conversely a congruence $f\equiv g_1\pmod r$ with $g_1\in L$ gives the integral form $(f-g_1)/r$, whose class has order $r$, so $r\mid r'$. Hence $r=r'=\#\T/(I+J)$. Now let $\phi_g\colon\T\to\mathcal{O}$ be the eigencharacter of $g$. It kills $J$ because $g\in f^\perp$, and $\phi_g\equiv\phi_f\pmod\lambda$ on all of $\T$ by hypothesis, so $\phi_g(I)\subset\lambda$. Thus $I+J$ lies in the kernel of $\T\to\mathcal{O}/\lambda$, a maximal ideal of residue characteristic $\ell$, and $\ell$ divides $\#\T/(I+J)=r$.
\end{proof}

\begin{lemma}[oldforms]\label{lem:old}
Let $h$ be a newform of weight $2$ on $\Gamma_0(M)$ with $M\mid N$ and $h\ne f$, with coefficients in $\mathcal{O}$, and let $\lambda\mid\ell$ be a prime of $\mathcal{O}$ such that $a_p(h)\equiv a_p(f)\pmod\lambda$ for all primes $p\nmid\ell N$. Assume that $p\nmid M$ for every prime $p\,\|\,N$ at which $\bar\rho$ is unramified. Then, after enlarging $\mathcal{O}$ and $\lambda$, the span of the forms $h(dz)$, $d\mid N/M$, contains an eigenform $g$ as in Lemma~\ref{lem:crit}. Consequently $\ell\mid r$.
\end{lemma}

\begin{proof}
The forms $h(dz)$, $d\mid N/M$, span a $\T$-stable subspace $V$ of $S_2(\Gamma_0(N))$ orthogonal to the newform $f$ (whether $M<N$, or $M=N$ and $h\ne f$). On $V$ the operators $T_n$ with $(n,N)=1$ act by the scalars $a_n(h)$, and $U_p$ for $p\mid M$ with $v_p(M)=v_p(N)$ acts by $a_p(h)$. For $p\mid N/M$ put $e_p=v_p(N/M)$; on the span of $h(p^jz)$, $0\le j\le e_p$, one has $U_ph(p^jz)=h(p^{j-1}z)$ for $j\ge1$, and $U_ph=a_p(h)h-p\,h(pz)$ if $p\nmid M$, $U_ph=a_p(h)h$ if $p\mid M$. The eigenvalues of $U_p$ there are therefore the roots $\alpha_p,\beta_p$ of $X^2-a_p(h)X+p$, together with $0$ if $e_p\ge2$, when $p\nmid M$; and $a_p(h)$, together with $0$ if $e_p\ge1$, when $p\mid M$. Since $V$ is the tensor product of these local spaces and $\T$ is commutative, for every choice of one such eigenvalue $\mu_p$ at each $p\mid N/M$ there is a common eigenform $g\in V\otimes\overline\Q$ of all of $\T$ with $U_pg=\mu_pg$ for $p\mid N/M$, $U_pg=a_p(h)g$ for the other $p\mid N$, and $T_ng=a_n(h)g$ for $(n,N)=1$; normalise it by $a_1(g)=1$. We choose the $\mu_p$ and check $a_p(g)\equiv a_p(f)\pmod\lambda$ prime by prime; the congruence for all $n$ then follows from the multiplicativity of the coefficients, from $a_{p^k}=a_p^k$ for $p\mid N$, and from the recursion $a_{p^{k+1}}=a_pa_{p^k}-p\,a_{p^{k-1}}$ for $p\nmid N$, which $f$ and $g$ both satisfy.

\emph{$p\nmid\ell N$.} $a_p(g)=a_p(h)\equiv a_p(f)$ by hypothesis. Consequently $\bar\rho_h\simeq\bar\rho$ by Chebotarev and Brauer--Nesbitt, $\bar\rho$ being irreducible.

\emph{$p=\ell$.} Both $f$ and $h$ have weight $2$ and level prime to $\ell$. By the theorems of Deligne and Fontaine describing $\bar\rho_f|_{D_\ell}$ and $\bar\rho_h|_{D_\ell}$ \cite[\S2]{Edixhoven92}, the residue class of $a_\ell$ is determined by the restriction to $D_\ell$: it is $0$ when that restriction is irreducible, and otherwise the eigenvalue of Frobenius on the unramified quotient. Hence $a_\ell(h)\equiv a_\ell(f)$.

\emph{$p^2\mid N$.} Here $E$ has additive reduction at $p$ and $a_p(f)=0$. If $v_p(M)=v_p(N)$ then $a_p(h)=0$, because a newform on $\Gamma_0(M)$ with $p^2\mid M$ has $a_p=0$ \cite[Theorem 3]{AtkinLehner70}. Otherwise $e_p\ge1$, and $e_p\ge2$ if $p\nmid M$, so $0$ is an available eigenvalue: take $\mu_p=0$.

\emph{$p\,\|\,N$, $\bar\rho$ unramified at $p$.} By hypothesis $p\nmid M$, so $e_p=1$ and $\mu_p\in\{\alpha_p,\beta_p\}$. Here $\rho_{E,\ell}|_{D_p}$ is an extension of an unramified character $\varepsilon$ by $\varepsilon\chi_\ell$ with $\varepsilon(\mathrm{Frob}_p)=a_p(f)=\pm1$ \cite[Theorem 3.1(e)]{DDT}, so the eigenvalues of $\bar\rho(\mathrm{Frob}_p)$ are $a_p(f)p$ and $a_p(f)$. As $\bar\rho_h\simeq\bar\rho$ is unramified at $p$ and $p\nmid M$, the characteristic polynomial of $\bar\rho_h(\mathrm{Frob}_p)$ is $X^2-a_p(h)X+p\equiv(X-a_p(f)p)(X-a_p(f))\pmod\lambda$, so one of $\alpha_p,\beta_p$ is congruent to $a_p(f)$ modulo a prime above $\lambda$ in $\mathcal{O}[\alpha_p]$; take it as $\mu_p$.

\emph{$p\,\|\,N$, $\bar\rho$ ramified at $p$.} Then $p\mid M$, since otherwise $\bar\rho_h$ would be unramified at $p$; so $v_p(M)=1=v_p(N)$ and $a_p(g)=a_p(h)=\pm1$. Both $\bar\rho|_{D_p}$ and $\bar\rho_h|_{D_p}$ are extensions of an unramified character by its twist by $\bar\chi_\ell$, with Frobenius eigenvalue $a_p(f)$, respectively $a_p(h)$, on the unramified quotient \cite[Theorem 3.1(e)]{DDT}, and both are non-split because $\bar\rho$ is ramified at $p$. A non-split extension of two characters has a unique stable line, so the isomorphism $\bar\rho|_{D_p}\simeq\bar\rho_h|_{D_p}$ matches the sublines and hence the quotients, and $a_p(h)\equiv a_p(f)$.

In all cases $a_p(g)\equiv a_p(f)\pmod\lambda$, and Lemma~\ref{lem:crit} applies.
\end{proof}

We use the optimal-level theorem of Ribet, Carayol and Diamond \cite{Ribet90,Carayol89,Diamond95} in the form of \cite[Theorem 3.15]{DDT}: \emph{let $\ell$ be odd and let $\bar\sigma\colon\Gal(\overline\Q/\Q)\to\mathrm{GL}_2(\overline{\F}_\ell)$ be irreducible and modular, absolutely irreducible on $\Gal(\overline\Q/\Q(\sqrt{-3}))$ if $\ell=3$; then $\bar\sigma$ arises from a newform of weight $2$ and level $N(\bar\sigma)\ell^{\delta(\bar\sigma)}$ with nebentypus of order prime to $\ell$, where $N(\bar\sigma)$ is the Serre conductor and $\delta(\bar\sigma)=0$ when $\bar\sigma|_{D_\ell}$ is good.} For $\bar\sigma=\bar\rho_{E',\ell}$ with $E'$ an elliptic curve of conductor prime to $\ell$, the restriction to $D_\ell$ comes from the finite flat group scheme $E'[\ell]$ and is good, and $\det\bar\sigma=\bar\chi_\ell$ forces the nebentypus, which is trivial modulo $\lambda$ and of order prime to $\ell$, to be trivial. So in that case $\bar\sigma$ arises from a newform of weight $2$ on $\Gamma_0(N(\bar\sigma))$. Recall also that $N(\bar\rho_{E',\ell})$ divides the conductor of $E'$ \cite{Carayol89}, and that $N(\bar\sigma)$ is prime to every prime at which $\bar\sigma$ is unramified.

\medskip\noindent\emph{Case (a).} Here $q^2\,\|\,N$ and $\ell=3$. By Lemma~\ref{lem:IV} the Serre conductor $M=N(\bar\rho)$ has exponent exactly $1$ at $q$, so $M\mid N/q$. The optimal-level theorem gives a newform $g$ of weight $2$ on $\Gamma_0(M)$ with $\bar\rho_g\simeq\bar\rho$, that is, $a_p(g)\equiv a_p(f)\pmod\lambda$ for all $p\nmid 3N$. Take $h=g$: it has level $M<N$, and $p\nmid M$ for every prime $p$ at which $\bar\rho$ is unramified. Lemma~\ref{lem:old} gives $3\mid r$.

\medskip\noindent\emph{Case (b).} Here $q^2\,\|\,N$, $\ell=3$, and $E$ has type $\mathrm{II}$ or $\mathrm{II}^*$ at $q\ge5$. Let $\chi=\chi_q$ and $E'=E^{(q)}$. By Lemma~\ref{lem:twist}, $E'$ has type $\mathrm{IV}^*$ or $\mathrm{IV}$ at $q$, the same types as $E$ elsewhere, and conductor $N$; its newform is $f'=f\otimes\chi$, a newform of weight $2$ on $\Gamma_0(N)$, and $\bar\rho_{E',3}=\bar\rho\otimes\chi$ satisfies the hypotheses of the optimal-level theorem since $\bar\rho$ does. By Lemma~\ref{lem:IV} applied to $E'$, the Serre conductor $M=N(\bar\rho\otimes\chi)$ has exponent exactly $1$ at $q$, and $M\mid N/q$. The optimal-level theorem gives a newform $g$ of weight $2$ on $\Gamma_0(M)$ with $a_p(g)\equiv a_p(f')=\chi(p)a_p(f)\pmod\lambda$ for all $p\nmid 3N$. As $q\,\|\,M$, Lemma~\ref{lem:twist}(iii) says that $h=g\otimes\chi$ is a newform on $\Gamma_0(Mq)$, and $Mq\mid N$. For $p\nmid 3N$,
\[
a_p(h)=\chi(p)a_p(g)\equiv\chi(p)^2a_p(f)=a_p(f)\pmod\lambda .
\]
Moreover $h\ne f$: otherwise $g=h\otimes\chi=f\otimes\chi=f'$ would have level $N$, whereas $M<N$. Finally, if $p\,\|\,N$ and $\bar\rho$ is unramified at $p$, then $p\ne q$, so $\bar\rho\otimes\chi$ is unramified at $p$ as well and $p\nmid M$, hence $p\nmid Mq$. Lemma~\ref{lem:old} gives $3\mid r$.

\medskip\noindent\emph{Case (c).} Here $E$ has type $\mathrm{I}_n^*$ at an odd prime $q\ne\ell$ with $\ell\mid n$, so $q^2\,\|\,N$. Let $\chi=\chi_q$ and $E'=E^{(q)}$. By Lemma~\ref{lem:twist}, $E'$ has multiplicative reduction of type $\mathrm{I}_n$ at $q$ and conductor $N/q$, so $f'=f\otimes\chi$ is a newform of weight $2$ on $\Gamma_0(N/q)$ with $q\,\|\,N/q$. By Lemma~\ref{lem:tamagawa-exp} applied to $E'$, $\bar\rho\otimes\chi=\bar\rho_{E',\ell}$ is unramified at $q$, so the Serre conductor $M=N(\bar\rho\otimes\chi)$ is prime to $q$ and $M\mid N/q^2$. The optimal-level theorem (this is the case of Ribet's theorem \cite{Ribet90}, $q$ exactly dividing the level and $\bar\rho\otimes\chi$ unramified at $q$) gives a newform $g$ of weight $2$ on $\Gamma_0(M)$ with $a_p(g)\equiv\chi(p)a_p(f)\pmod\lambda$ for all $p\nmid\ell N$. As $q\nmid M$, $h=g\otimes\chi$ is a newform on $\Gamma_0(Mq^2)$ with $Mq^2\mid N$, congruent to $f$ at all $p\nmid\ell N$ by the computation of case (b), and $h\ne f$ since $g=h\otimes\chi$ would otherwise equal $f'$, of level $N/q\ne M$. The hypothesis of Lemma~\ref{lem:old} holds as in case (b), and $\ell\mid r$. \qed

\begin{example}
The congruences produced by the proof can be seen directly on the tables (labels as in \cite{Cremona97}). (b): the curve 121a1 has type $\mathrm{II}$ at $11$; its twist by $-11$ lies in the class 121c and has type $\mathrm{IV}^*$ at $11$; its mod $3$ representation has conductor $11$ and is congruent to the form of 11a1; twisting back, $f_{\text{11a}}\otimes\chi_{-11}$ is the form of 121d, and indeed $a_p(\text{121a1})\equiv\chi_{-11}(p)\,a_p(\text{11a1})\pmod 3$ for all $p\ne3,11$ (checked for $p<500$), while $\deg\varphi_{\text{121a1}}=6$: here $e_3(E,11)=v_3(11+1)=1$ and the bound is attained. (c) at $\ell=3$: 539c1 has type $\mathrm{I}_3^*$ at $7$; its twist by $-7$ has conductor $77$ and type $\mathrm{I}_3$ at $7$, and $a_p(\text{539c1})\equiv\chi_{-7}(p)\,a_p(\text{11a1})\pmod3$ for all $p\ne3,7,11$ (checked for $p<1000$); $\deg\varphi_{\text{539c1}}=2^5\cdot3^2$, and the two terms $e_3(E,7)=v_3(48)=1$ and $t_3(E,7)=v_3(3)=1$ are both attained. (c) at $\ell=5$: 980i1 has type $\mathrm{I}_5^*$ at $7$; its twist by $-7$ has conductor $140$ and type $\mathrm{I}_5$ at $7$, and $a_p(\text{980i1})\equiv\chi_{-7}(p)\,a_p(\text{20a1})\pmod5$ for all $p\ne2,5,7$ (checked for $p<1000$); $\deg\varphi_{\text{980i1}}=2^6\cdot3^2\cdot5$.
\end{example}

\begin{remark}
The hypothesis $\ell\nmid N$ is used in the optimal-level theorem, in Lemma~\ref{lem:old} at $p=\ell$, and to pass from the congruence number to the modular degree. The last step needs only $\ell^2\nmid N$, and it cannot be dropped there: Table 1 of \cite{ARS12} lists the curve 99A1, of conductor $9\cdot11$, with modular degree $4$ and congruence number $12$. The data nevertheless show the conclusion of the theorem for every curve with $\ell\,\|\,N$ or $\ell^2\mid N$ and a prime of one of the three kinds; so do they at $q=3$ in cases (a) and (b), where Lemma~\ref{lem:IV} does not apply because the ramification is wild. The proof gives one factor of $\ell$ for each such prime but says nothing about additivity with the multiplicative contributions or with each other; for the types $\mathrm{I}_n^*$ and $\mathrm{I}_0^*$ that is done in the next section, where the proof of case (c) is read quantitatively.
\end{remark}

\section{The quantitative theorems}\label{sec:quant}

\subsection{Congruence ideals and the formula of Diamond, Flach and Guo}\label{sec:dfg}
Let $g$ be a newform of weight $2$ on $\Gamma_0(M)$ with rational coefficients, $\lambda=\ell$ a prime not dividing $2M$, and $\Sigma$ a finite set of primes. Put $M^\Sigma=M\prod_{p\in\Sigma}p^{\delta_p}$ and let $g^\Sigma=\sum_{(n,\Sigma)=1}a_n(g)q^n$ be the $\Sigma$-depleted eigenform, of level $M^\Sigma$. The integral homology $H_1(X_0(M^\Sigma),\Z_\ell)$ carries the intersection pairing, perfect and alternating; composed with the Atkin--Lehner involution $w_{M^\Sigma}$ it becomes a perfect alternating pairing for which the full Hecke algebra, the operators $U_p$ included, is self-adjoint \cite[\S4.4]{DDT}. Its restriction to the rank-two lattice $M_{g^\Sigma,\lambda}$ killed by the kernel of the eigencharacter of $g^\Sigma$ has a Gram determinant, which is a square since the pairing is alternating, and \cite[\S1.7.3]{DFG04} define the \emph{naive $\Sigma$-finite congruence ideal} $\eta^\Sigma_{g,\lambda}\subset\Z_\ell$ as its square root, an integral ideal.

\begin{theorem}[{\cite[Proposition 1.4(c)]{DFG04}}]\label{thm:dfg14}
If $\bar\rho_g$ is irreducible then
\[
\eta^\Sigma_{g,\lambda}=\eta^{\emptyset}_{g,\lambda}\prod_{p\in\Sigma}L_p^{\nv}(A_g,1)^{-1}.
\]
\end{theorem}
Only the primes with $\delta_p\ge1$ contribute: adding an additive prime to $\Sigma$ changes neither the level nor the ideal; adding a good prime $p$ multiplies the ideal by $(1-\frac\alpha\beta p^{-1})(1-p^{-1})(1-\frac\beta\alpha p^{-1})$, Wiles's level-raising factor; adding a multiplicative prime multiplies it by $1-p^{-2}$.

We also need the following piece of commutative algebra, which is the inequality half of \cite[Lemma 4.17]{DDT}. Let $\T$ be a commutative $\mathcal{O}$-algebra, finite and free as an $\mathcal{O}$-module, acting on an $\mathcal{O}$-module $L$ free of finite rank with a perfect $\mathcal{O}$-bilinear pairing for which $\T$ is self-adjoint; let $\pi\colon\T\to\mathcal{O}$ be a surjective $\mathcal{O}$-algebra homomorphism, $\wp=\ker\pi$, $I=\operatorname{Ann}_\T(\wp)$ and $\eta=\pi(I)$, the congruence ideal of $\pi$; and suppose that $L[\wp]$ is free of rank $d$ over $\mathcal{O}$ with Gram determinant $\Delta\ne0$.

\begin{lemma}\label{lem:ddt417}
$\eta^d\subseteq\Delta\mathcal{O}$, that is, $d\cdot v(\eta)\ge v(\Delta)$; and equality holds if $L$ is free over $\T$.
\end{lemma}
\begin{proof}
Since $L$ is torsion-free and $\wp\cap I=0$, the sum $L[\wp]+L[I]$ is direct. For $t\in I$ with $\pi(t)$ a generator of $\eta$ and $x\in L$, one has $tx\in L[\wp]$ (as $\wp I=0$) and $(\pi(t)-t)x\in L[I]$ (as $s(\pi(t)-t)\in\wp\cap I=0$ for $s\in I$), so $Q=L/(L[\wp]\oplus L[I])$ is killed by $\eta$. The pairing identifies $L/L[\wp]^\perp$ with $\operatorname{Hom}(L[\wp],\mathcal{O})$, because $L[\wp]$ is saturated, and $L[\wp]^\perp=L[I]$: for $y\in L[I]$ and $x\in L[\wp]$, $\pi(t)\langle x,y\rangle=\langle tx,y\rangle=\langle x,ty\rangle=0$; conversely if $y\perp L[\wp]$ then $\langle ty,z\rangle=\langle y,tz\rangle=0$ for all $z$, since $tz\in L[\wp]$, so $ty=0$. Hence $Q$ is the cokernel of $L[\wp]\to\operatorname{Hom}(L[\wp],\mathcal{O})$, of order $\#\mathcal{O}/\Delta$, and it is generated by $d$ elements and killed by $\eta$, so $\#\mathcal{O}/\Delta\le\#(\mathcal{O}/\eta)^d$. If $L\simeq\T^d$ then $Q\simeq(\T/(\wp+I))^d\simeq(\mathcal{O}/\eta)^d$.
\end{proof}

\subsection{Twist monotonicity of the congruence number}
\begin{theorem}\label{thm:twist}
Let $E/\Q$ have conductor $N$, $S$ a set of odd primes, $D=\prod_{q\in S}q$, $\chi_S$ the product of the quadratic characters of conductor $q\in S$, and $E_S=E\otimes\chi_S$ of conductor $N_S$. If $q^2\mid N$ for every $q\in S$ and $\operatorname{lcm}(N_S,D^2)\mid N$, then $r_{E_S}\mid r_E$; if moreover $\operatorname{lcm}(N,D^2)\mid N_S$, then $r_{E_S}=r_E$.
\end{theorem}
\begin{proof}
Let $f=f_E$, $g=f_{E_S}$ and $\iota\colon h\mapsto h\otimes\chi_S=\sum\chi_S(n)a_n(h)q^n$, which maps $S_2(\Gamma_0(N_S))$ into $S_2(\Gamma_0(\operatorname{lcm}(N_S,D^2)))\subset S_2(\Gamma_0(N))$, preserves integrality, and sends an eigenform of all $T_n$, $(n,N_S)=1$, to an eigenform with eigenvalues twisted by $\chi_S(n)$. Then $\iota(g)=f$: both are multiplicative, $a_p(f)=\chi_S(p)a_p(g)$ for $p\nmid D$ since the twist is unramified outside $S$, and $a_{q^k}(f)=0$ for $q\in S$ since $q^2\mid N$. If $h$ is an eigenform of level dividing $N_S$ whose $T_p$-eigenvalues differ from those of $g$ at infinitely many $p$, the same is true of $\iota(h)$ and $f$, so $\iota(h)\in f^\perp$; hence $\iota$ maps $g^\perp$, the $\T$-stable complement of the $g$-line (the sum of the generalised eigenspaces of the other systems of eigenvalues), into $f^\perp$. A congruence $g\equiv h\pmod r$ with $h\in g^\perp$ integral therefore gives $f\equiv\iota(h)\pmod r$ with $\iota(h)\in f^\perp$ integral, and $r\mid r_E$ by Lemma~\ref{lem:crit}. With $r=r_{E_S}$ this is the first claim; the second hypothesis makes the situation symmetric.
\end{proof}
The hypotheses hold when (a) $E$ has type $\mathrm{I}_n^*$ at each $q\in S$ ($N_S=N/D$), (b) type $\mathrm{I}_0^*$ at each $q\in S$ ($N_S=N/D^2$, $E_S$ good at $S$), (c) type $\mathrm{II}$, $\mathrm{III}$, $\mathrm{IV}$ or a dual at each $q\in S$, $q\ge5$ ($N_S=N$, equality). On the working set, for the $157{,}938$ pairs (curve, prime of type $\mathrm{I}_n^*$) and the $292{,}580$ pairs (curve, prime of potentially good additive type), $v_\ell(\degphi)\ge v_\ell(\deg\varphi_{E_q})$ holds for $\ell\in\{3,5,7,11,13\}$, $\ell^2\nmid N$, with no exception, with equality in every pair of case (c) and strict inequality only in case (b), where the difference is the Euler term of Theorem~\ref{thm:euler}.

\subsection{The Euler terms of the types \texorpdfstring{$\mathrm{I}_0^*$ and $\mathrm{I}_n^*$}{I0* and In*}}

\begin{theorem}\label{thm:euler}
Let $\ell$ be an odd prime with $\ell\nmid 2N$ and $E[\ell]$ irreducible. Let $S$ be a set of primes $q\ge5$ at which $E$ has type $\mathrm{I}_0^*$ or $\mathrm{I}_n^*$, and $E_S=E\otimes\chi_S$. Then
\[
v_\ell(\degphi)\ \ge\ v_\ell(r_{E_S})\ +\ \sum_{q\in S}e_\ell(E,q),
\]
with $e_\ell(E,q)=v_\ell(q-1)+v_\ell((q+1)^2-a_q(E_S)^2)$ for type $\mathrm{I}_0^*$ and $v_\ell(q^2-1)$ for type $\mathrm{I}_n^*$.
\end{theorem}

\begin{proof}
Write $f=f_E$, $g=f_{E_S}$ of level $N_S$, $D=\prod_{q\in S}q$, $\Sigma=S$, and $\mathcal{O}=\Z_\ell$. By Lemma~\ref{lem:twist}, $E_S$ has good reduction at the $\mathrm{I}_0^*$ primes of $S$ and multiplicative reduction at the $\mathrm{I}_n^*$ primes, so $\delta_q=2$, resp.\ $1$, for $g$ at $q\in S$, and $N_S^\Sigma=N_S\prod_{q\in S}q^{\delta_q}=N$: the $S$-depleted form $g^S$ lives at level $N$, and $\iota(g^S)=f$ with $\iota$ the twist map of Theorem~\ref{thm:twist}, which kills the coefficients at multiples of $D$ and so does not distinguish $g$ from $g^S$.

\emph{Step 1: the twist map and the full Hecke algebra.} Let $\T$ be the full Hecke algebra of level $N$, $\pi\colon\T\to\Z$ the eigencharacter of $g^S$, and $\mathfrak{m}$ the maximal ideal of $\T_\ell=\T\otimes\Z_\ell$ containing $\ker\pi$ and $\ell$; it contains $U_q$ for $q\in S$. The congruence ideal of $\pi$, $\eta^{\mathrm{Hecke}}(g^S)=\pi(\operatorname{Ann}_\T\ker\pi)$, is computed in the localisation $\T_{\mathfrak m}$, since the other local factors of $\T_\ell$ are killed by $\pi$ and lie in the annihilator; and the generalised eigenspace of $\pi$ in $S_2(\Gamma_0(N))\otimes\overline\Q$ is the line of $g^S$ (strong multiplicity one identifies the forms with the $T_p$-eigenvalues of $g$, $p\nmid N$, with the span of the $g(dz)$, $d\mid N/N_S$, and on that span the operators $U_q$, $q\in S$, have the eigenvalue $0$ with multiplicity one). The eigenforms of $\T$ whose eigencharacter is congruent to $\pi$ modulo $\mathfrak m$ have $U_q\equiv0$ at $q\in S$. Among the stabilisations of $g$, only $g^S$ qualifies: the others have $U_q$-eigenvalue $a_q(g)=\pm1$ (multiplicative case) or a root $\alpha,\beta$ of $X^2-a_q(g)X+q$ (good case), which are $\ell$-adic units since $\ell\ne q$. So the $\mathfrak m$-component of $S_2(\Gamma_0(N);\Z_\ell)$ is spanned by $g^S$ and by eigenforms $h'$ with $U_qh'=0$ for $q\in S$ (the available eigenvalues being $\alpha,\beta,0$ or $\pm1,0$ or $0$, by the proof of Lemma~\ref{lem:old}) arising from newforms $h\ne g$ of level dividing $N$. For such an $h'$, $\iota(h')=h'\otimes\chi_S$ is again an eigenform of $\T$ at level $N$ (its coefficients at multiples of $D$ vanish, $U_p$ for $p\nmid D$ commutes with $\iota$ up to the scalar $\chi_S(p)$, and the level of a twist by $\chi_S$ divides $\operatorname{lcm}(N,D^2)=N$), whose $T_p$-eigenvalues $\chi_S(p)a_p(h)$, $p\nmid DN$, differ from those of $f$ at infinitely many $p$; so $\iota(h')\in f^\perp$. Now the proof of Lemma~\ref{lem:crit}, carried out over $\Z_\ell$ with the pairing $S_2(\Gamma_0(N);\Z_\ell)\times\T_\ell\to\Z_\ell$, shows that $\ell^{v_\ell(\eta^{\mathrm{Hecke}}(g^S))}$ is the largest power $\ell^t$ for which $g^S\equiv h\pmod{\ell^t}$ coefficientwise with $h\in S_2(\Gamma_0(N);\Z_\ell)$ in the $\T$-stable complement of the $g^S$-line; and projecting $h$ to the $\mathfrak m$-component (an idempotent of $\T_\ell$, so the projection stays in the complement and the congruence survives) we may take $h$ to be a $\overline{\Z}_\ell$-combination of the eigenforms $h'$ above. Then $f=\iota(g^S)\equiv\iota(h)\pmod{\ell^t}$ with $\iota(h)\in f^\perp$ integral, and the $\Z_\ell$-version of Lemma~\ref{lem:crit} for $f$ gives
\[
v_\ell(r_E)\ \ge\ v_\ell(\eta^{\mathrm{Hecke}}(g^S)).
\]

\emph{Step 2: Hecke versus cohomological congruence ideal.} Apply Lemma~\ref{lem:ddt417} to the module $L=H_1(X_0(N),\Z_\ell)_{\mathfrak m}$, with the self-adjoint perfect alternating pairing of Section~\ref{sec:dfg}, and to $\pi$: $L[\wp]$ is the rank-two lattice $M_{g^S,\lambda}$, whose Gram determinant is $(\eta^S_{g,\lambda})^2$ up to a unit. Hence
\[
v_\ell(\eta^{\mathrm{Hecke}}(g^S))\ \ge\ v_\lambda(\eta^S_{g,\lambda}).
\]
Freeness of $L$ over $\T_{\mathfrak m}$ is not needed for this inequality.

\emph{Step 3: the formula of Diamond, Flach and Guo.} By Theorem~\ref{thm:dfg14} applied to $g$ (its residual representation $\bar\rho\otimes\chi_S$ is irreducible and $\ell\nmid 2N_S$),
\[
v_\lambda(\eta^S_{g,\lambda})=v_\lambda(\eta^\emptyset_{g,\lambda})+\sum_{q\in S}v_\ell\bigl(L_q^{\nv}(A_g,1)^{-1}\bigr)=v_\lambda(\eta^\emptyset_{g,\lambda})+\sum_{q\in S}e_\ell(E,q),
\]
the last equality because $A_g=A_f$ and, at $q\in S$, $g$ has $\delta_q\ge1$, so that its naive factor is the true one, which is the factor that $f$ (with $a_q(f)=0$) omits; the values are those of Lemma~\ref{lem:euler}.

\emph{Step 4: the level $N_S$.} At level $N_S$ the form $g$ is new, $\ell\nmid N_S$ and $\bar\rho_g$ is irreducible with Serre weight $2\ne\ell$, so multiplicity one holds for the maximal ideal of $g$ (\cite[Theorem 3.5]{RibetStein} and the references in its proof; the only possible exception there has Serre weight $\ell$), hence the localised Hecke algebra is Gorenstein and $H_1(X_0(N_S),\Z_\ell)_{\mathfrak m}$ is free of rank two over it \cite{Tilouine97}; the equality case of Lemma~\ref{lem:ddt417} then gives $v_\lambda(\eta^\emptyset_{g,\lambda})=v_\ell(\eta^{\mathrm{Hecke}}(g))=v_\ell(r_{E_S})$, the last by Lemma~\ref{lem:crit}. Putting the steps together and using $v_\ell(\degphi)=v_\ell(r_E)$ (as $\ell\nmid N$) proves the theorem. Without Step 4 one still has $v_\ell(\degphi)\ge\sum_{q\in S}e_\ell(E,q)$, since $\eta^\emptyset_{g,\lambda}$ is an integral ideal.
\end{proof}

\subsection{The Tamagawa terms}
Kim and Ota \cite[Theorem 1.1]{KimOta} prove, for a newform $f$ of weight $2$ and level $N=N^+N^-$ with $N^-$ square-free, a prime $\ell\ge5$ with $\ell\nmid N$ and $\bar\rho_f$ absolutely irreducible on $\Gal(\overline\Q/\Q(\sqrt{\ell^*}))$, and under the condition that $\bar\rho_f$ be ramified at every $q\mid N^-$ with $q\equiv\pm1\pmod\ell$,
\begin{equation}\label{eq:ko}
\ord_\lambda\eta_f(N)=\ord_\lambda\eta_f(N^+,N^-)+\sum_{q\mid N^-}t_f(q),
\end{equation}
where $\eta_f(N)$ is the congruence ideal of $f$ in the full Hecke algebra of level $N$, $\eta_f(N^+,N^-)$ the congruence ideal in its $N^-$-new quotient, an integral ideal, and $t_f(q)$ the Tamagawa exponent \cite[Definition 2.12]{KimOta}, which is $t_\ell(E,q)$ for an elliptic curve by Lemma~\ref{lem:tamagawa-exp}.

\begin{theorem}\label{thm:C}
Let $\ell\ge5$, $\ell\nmid N$, with $\bar\rho_{E,\ell}$ absolutely irreducible on $\Gal(\overline\Q/\Q(\sqrt{\ell^*}))$. Let $M$ be the set of primes $q$ at which $E$ has type $\mathrm{I}_n$, or $\mathrm{I}_n^*$ with $q$ odd, with $\ell\mid n$, and $M'\subset M$ the subset of primes $q\not\equiv\pm1\pmod\ell$. Let $S$ be the set of $\mathrm{I}_n^*$ primes in $M'$. Then $v_\ell(r_{E_S})\ge\sum_{q\in M'}t_\ell(E,q)$, and hence $v_\ell(\degphi)\ge\sum_{q\in M'}t_\ell(E,q)$.
\end{theorem}
\begin{proof}
$E_S$ is multiplicative at every prime of $M'$ with the same index (Lemma~\ref{lem:twist}), has conductor prime to $\ell$, and $\bar\rho_{E_S}=\bar\rho\otimes\chi_S$ satisfies the irreducibility hypothesis. Apply \eqref{eq:ko} to $f_{E_S}$ with $N^-=\prod_{q\in M'}q$, square-free and exactly dividing the level; the ramification condition holds vacuously. Since $\eta_f(N^+,N^-)$ is integral, $v_\ell(r_{E_S})\ge\sum_{M'}t_\ell$, and Theorem~\ref{thm:twist}(a) gives $r_{E_S}\mid r_E$, with $v_\ell(r_E)=v_\ell(\degphi)$.
\end{proof}

\begin{corollary}\label{cor:both}
Under the hypotheses of Theorem~\ref{thm:C}, let $S$ be the set of primes $q\ge5$ of type $\mathrm{I}_0^*$ or $\mathrm{I}_n^*$. Then
\[
v_\ell(\degphi)\ \ge\ \sum_{q\in S}e_\ell(E,q)+\sum_{q\in M'}t_\ell(E,q).
\]
\end{corollary}
\begin{proof}
Theorem~\ref{thm:euler} for $S$ gives $v_\ell(\degphi)\ge v_\ell(r_{E_S})+\sum_S e_\ell$. The twist $E_S$ is multiplicative at the $\mathrm{I}_n^*$ primes of $S$ and has the same types as $E$ elsewhere, so its set $M'$ is that of $E$ with the same indices, and Theorem~\ref{thm:C} applied to $E_S$ (which twists further at the $\mathrm{I}_n^*$ primes of $M'$ not in $S$, that is, at $q=3$) gives $v_\ell(r_{E_S})\ge\sum_{M'}t_\ell$.
\end{proof}

\begin{proof}[Proof of Theorem~\ref{thm:quant}]
The first inequality is Theorem~\ref{thm:euler} with the integrality of $r_{E_S}$; the second is Corollary~\ref{cor:both}.
\end{proof}

\begin{remark}
(1) Theorem~\ref{thm:euler} is the quantitative form of the twist-and-lower argument of Section~\ref{sec:proof}, read in the other direction: the level is raised, from $N_S$ to $N_S\prod q^{\delta_q}=N$, and the price is the level-raising Euler factor. For $\ell=3$ it says that a single prime of type $\mathrm{I}_n^*$, with $3\mid n$ or not, forces $3\mid\degphi$, and that a prime of type $\mathrm{I}_0^*$ with $q\equiv1\pmod3$ does too; neither is in Theorem~\ref{thm:qual}. (2) The identification of the pairing of \cite[\S1.7.3]{DFG04} with the one used in Step 2 is up to units and is the point at which the citation of \cite{DFG04} should be checked by the reader. (3) Step 1 is where the argument fails for the potentially good types $\mathrm{II}$, $\mathrm{IV}$, $\mathrm{III}$: there the twist does not lower the level, $f$ is its own depletion, and Theorem~\ref{thm:dfg14} gives no factor; the omitted factor sits inside the naive $L$-value, which is where the next section picks it up.
\end{remark}

\section{The remaining types: a Selmer-group reformulation}\label{sec:selmer}

Let $\Sigma$ be the set of primes dividing $N$, $\mathcal{O}=\Z_\ell$, $K=\Q_\ell$ and $M=\ad V\otimes K/\mathcal{O}$. For $\lambda\nmid2N$ with $\bar\rho$ absolutely irreducible on $\Gal(\overline\Q/\Q(\sqrt{\ell^*}))$, \cite[Theorem 0.2]{DFG04} (their Theorem 3.7) says that the $\Sigma$-imprimitive adjoint Selmer group $H^1_\Sigma(\Q,M)$, with no condition at the primes of $\Sigma$ and the Bloch--Kato condition elsewhere, is finite of length $v_\ell(\eta_f^\Sigma)$; and Theorem~\ref{thm:dfg14}, with multiplicity one at level $N$ identifying $\eta^\emptyset_f$ with $r_E$ as in Step 4 above, gives $v_\ell(\eta_f^\Sigma)=v_\ell(r_E)+\sum_{p\,\|\,N}v_\ell(p^2-1)$. So the Euler-factor rule, for all types at once, is the statement
\begin{equation}\label{eq:selmer}
\operatorname{length}H^1_\Sigma(\Q,M)\ \ge\ \sum_{p\in\Sigma,\ \delta_p=0}e_\ell(E,p)\ +\ \sum_{p\in\Sigma,\ \delta_p=1}\bigl(t_\ell(E,p)+v_\ell(p^2-1)\bigr).
\end{equation}
Compare $H^1_\Sigma$ with the Bloch--Kato Selmer group $H^1_f(\Q,M)\subseteq H^1_\Sigma(\Q,M)$ through the Greenberg--Wiles formula \cite[Theorem 2.18]{DDT} applied to the finite layers $M[\ell^n]$ with the two sets of local conditions. The local index at $p\in\Sigma$ between the unrestricted and the unramified condition is, by local Tate duality, $\#H^0(\Q_p,M[\ell^n]^*)$, and in the limit, when $\ell$ does not divide the order of $\rho(I_p)$ and $H^0(\Q_p,\ad\rho(1)\otimes K/\mathcal{O})$ is finite, this is $\#\mathcal{O}/\det(1-\Frob_p\mid(\ad V)^{I_p}(1))$, the $\ell$-part of $L_p(A_f,1)^{-1}$: $\ell^{e_\ell(E,p)}$ at an additive prime and $\ell^{v_\ell(p^2-1)}$ at a multiplicative one. The global correction is the ratio of the dual Selmer groups, $\#H^1_{f^*}(\Q,M^*)/\#H^1_{\Sigma^*}(\Q,M^*)\ge1$, with $M^*=\ad\rho(1)\otimes K/\mathcal{O}$, whose Bloch--Kato conjecture \cite{DFG04} also prove. Hence, under these hypotheses, \eqref{eq:selmer} is equivalent to:

\begin{conjecture}[Selmer form of the rule]\label{conj:selmer}
With $\ell\nmid2N$ and $\bar\rho$ absolutely irreducible on $\Gal(\overline\Q/\Q(\sqrt{\ell^*}))$,
\begin{multline*}
\operatorname{length}H^1_f(\Q,\ad\rho\otimes K/\mathcal{O})\ +\ \sum_{p\in\Sigma}\operatorname{length}H^0(\Q_p,\ad\rho(1)\otimes K/\mathcal{O})\\
-\ \operatorname{length}\frac{H^1_{f^*}(\Q,\ad\rho(1)\otimes K/\mathcal{O})}{H^1_{\Sigma^*}(\Q,\ad\rho(1)\otimes K/\mathcal{O})}\ \ \ge\ \ \text{right-hand side of \eqref{eq:selmer}}.
\end{multline*}
\end{conjecture}

The data of Section~\ref{sec:data} say more: the Bloch--Kato Selmer group supplies the Tamagawa exponents, the local terms supply the Euler terms, and the dual correction vanishes, except in the type $\mathrm{III}$ configuration of Section~\ref{sec:anomaly}, where it equals one. The exact local factors are thus the $H^0(\Q_p,\ad\rho(1)\otimes K/\mathcal{O})$, whose residual versions $H^0(\Q_p,\ad\bar\rho(1))$ are the local terms in the Taylor--Wiles count of the tangent space \cite[Corollary 2.43]{DDT}; the twist-symmetrised weights of the earlier version were their shadow, explained by the twist invariance of $\ad\rho$. The one thing we do not find in the literature is the vanishing of the dual correction, whose failure the type $\mathrm{III}$ anomaly exhibits; the refined integral structures of \cite[\S1.8]{DFG04}, built to account for the Euler factors missing from the naive $L$-function at the exceptional primes, are the natural place to settle it, and the Taylor--Wiles machinery of \cite[\S2]{KimOta} controls the same groups.

\section{The computations}\label{sec:data}

\subsection{Data and protocol}
We used Cremona's tables \cite{Cremona97} in the form of the \texttt{ecdata} repository at commit \texttt{25cec5e}, which contain every elliptic curve over $\Q$ of conductor below $500000$ with its Kodaira symbols, Tamagawa numbers, modular degree, isogeny class, torsion and, up to conductor $400000$, the images of the mod $\ell$ Galois representations. Optimality is determined in the tables for every class of conductor at most $400000$ and for $359{,}009$ further classes up to $500000$. The ranges are the working set ($N\le300000$), the hold-out ($300000<N\le400000$), the extension ($400000<N\le500000$) and the out-of-table set of Section~\ref{sec:oot}. The component-group inequality \eqref{eq:main} and Conjecture~\ref{conj:eis} were found on the working set, written down, and then tested once on the hold-out, which had been held back, and afterwards on the extension; nothing was adjusted after the hold-out test. The twist-symmetrised form \eqref{eq:twist} replaced an earlier, weaker formulation of the same phenomenon on the working set and was then tested once on the hold-out and the extension. Conjecture~\ref{conj:euler} came last: it was formulated on the working set from the minima by Kodaira type (Section~\ref{sec:iso}), the terms at $2$ and $3$ were fixed afterwards by PARI's Euler factors, and the rule was then tested on all four ranges. Since the hold-out and the extension had already been used for \eqref{eq:main} and \eqref{eq:twist}, they are a confirmation and not a blind test for the rule; the out-of-table curves, generated before the rule was formulated, are. The Kodaira types and discriminant exponents at the bad primes were recomputed with PARI/GP 2.15 \cite{PARI} from the $a$-invariants, the Euler factors of $\Sym f$ at every bad prime were computed with \texttt{lfunsympow} and \texttt{lfuneuler}, and fifteen modular degrees spanning the configurations were recomputed with \texttt{ellmoddegree}, all agreeing with the tables.

\subsection{The rule on whole ranges}
Table~\ref{tab:uniform} checks \eqref{eq:rule}, with PARI's Euler factors at every bad prime $q\ne\ell$ and the terms $t_\ell$, $w_3$ as defined, on the four ranges at $\ell=3,5,7,11,13$: no violation at any $\ell$. It also reports what happens when the type $\mathrm{III}$ terms are added at every prime (exceptions at $\ell=3$ only, Section~\ref{sec:anomaly}) and when the weights at $q=\ell$ are dropped (no violation, so the weights $w_3$ only sharpen the bound; at $\ell\ge5$ nothing is needed at $q=\ell$).

\begin{table}[ht]\centering\footnotesize
\begin{tabular}{llrrrrrr}
\toprule
set & $\ell$ & curves & \eqref{eq:rule} $>0$ & viol. & equality & viol.\ $+\mathrm{III}$ & no $q=\ell$: $>0$ / viol. / equality\\
\midrule
working set & 3 & 1,233,729 & 861,074 & 0 & 56.8\% & 364 & 800,833 / 0 / 52.1\% \\
 & 5 & 1,303,224 & 330,477 & 0 & 75.5\% & 0 & 330,477 / 0 / 75.5\% \\
 & 7 & 1,307,753 & 173,488 & 0 & 84.1\% & 0 & 173,488 / 0 / 84.1\% \\
 & 11 & 1,309,103 & 59,495 & 0 & 91.5\% & 0 & 59,495 / 0 / 91.5\% \\
 & 13 & 1,308,945 & 35,635 & 0 & 92.4\% & 0 & 35,635 / 0 / 92.4\% \\
\midrule
hold-out & 3 & 409,577 & 295,973 & 0 & 53.9\% & 68 & 277,265 / 0 / 49.3\% \\
 & 5 & 428,251 & 114,860 & 0 & 74.0\% & 0 & 114,860 / 0 / 74.0\% \\
 & 7 & 429,280 & 60,580 & 0 & 83.2\% & 0 & 60,580 / 0 / 83.2\% \\
 & 11 & 429,502 & 21,097 & 0 & 90.6\% & 0 & 21,097 / 0 / 90.6\% \\
 & 13 & 429,498 & 12,681 & 0 & 92.3\% & 0 & 12,681 / 0 / 92.3\% \\
\midrule
extension & 3 & 348,227 & 255,507 & 0 & 53.2\% & 43 & 238,359 / 0 / 48.1\% \\
 & 5 & 358,244 & 99,401 & 0 & 73.2\% & 0 & 99,401 / 0 / 73.2\% \\
 & 7 & 358,717 & 52,766 & 0 & 82.7\% & 0 & 52,766 / 0 / 82.7\% \\
 & 11 & 358,830 & 19,136 & 0 & 90.6\% & 0 & 19,136 / 0 / 90.6\% \\
 & 13 & 358,816 & 11,484 & 0 & 91.9\% & 0 & 11,484 / 0 / 91.9\% \\
\midrule
out of table & 3 & 893 & 827 & 0 & 51.6\% & 0 & 804 / 0 / 42.3\% \\
 & 5 & 893 & 111 & 0 & 82.0\% & 0 & 111 / 0 / 82.0\% \\
 & 7 & 893 & 52 & 0 & 90.4\% & 0 & 52 / 0 / 90.4\% \\
 & 11 & 893 & 19 & 0 & 94.7\% & 0 & 19 / 0 / 94.7\% \\
 & 13 & 893 & 12 & 0 & 100.0\% & 0 & 12 / 0 / 100.0\% \\
\bottomrule
\end{tabular}

\caption{The Euler-factor rule with PARI's factors at all primes: curves with $E[\ell]$ irreducible, not CM, with a positive bound in \eqref{eq:rule}; violations; equality rate among them; violations when the $\mathrm{III}$ terms are added at every prime (``$+\mathrm{III}$''); and the same quantities without the weights at $q=\ell$ (``no $q=\ell$'').}\label{tab:uniform}
\end{table}

\subsection{Calibration and sharpness}\label{sec:sharp}
The zero counts are not tautological. Among the $372{,}655$ working-set curves with $E[3]$ irreducible, not CM, and right-hand side $0$ in \eqref{eq:rule}, the prime $3$ does not divide $\degphi$ for $52.3\%$; the corresponding proportions are $72.7\%$, $82.9\%$, $90.4\%$ and $92.2\%$ at $\ell=5,7,11,13$. If the valuation of the degree were independent of the local data, the $861{,}074$ working-set curves with a positive bound at $\ell=3$ would produce roughly $590{,}000$ violations; there are none. The bounds are attained: among working-set curves with a positive bound, $v_\ell(\degphi)$ equals the right-hand side of \eqref{eq:rule} in $56.8\%$ of the cases at $\ell=3$, $75.5\%$ at $\ell=5$, $84.1\%$ at $\ell=7$, $91.5\%$ at $\ell=11$ and $92.4\%$ at $\ell=13$, and the mean excess at $\ell=3$ is $0.64$. The whole odd part of $\degphi$ is another matter: for the $1{,}226{,}216$ working-set optimal non-CM curves with no isogeny of odd degree, the odd part of $\degphi$ equals $\prod_\ell\ell^{\,b_\ell}$, with $b_\ell$ the right-hand side of \eqref{eq:twist}, in only $1.85\%$ of cases; the quotient is most often $9$, $3$, $15$, $45$, $27$ or $21$, its smallest prime factor is $3$ in $63\%$ of the cases and divides $N$ in $52\%$. So the local factors account for a definite, sharp part of the modular degree and for little of the rest, which is governed by congruences of other origins and by the size of the Petersson norm.

\subsection{Each term in isolation}\label{sec:iso}
For working-set curves with exactly one additive prime $q\ge5$, $q\ne\ell$, and no multiplicative prime with $\ell\mid n$, Table~\ref{tab:iso} gives the minimum of $v_\ell(\degphi)$ by Kodaira type and by $e_\ell(E,q)$. In every cell, for $\ell=3,5,7$, the minimum equals $e_\ell$ and every curve has $v_\ell(\degphi)\ge e_\ell$; for $\mathrm{I}_n^*$ the excess over $e_\ell$ has minimum exactly $v_\ell(n)$ in every cell. The cells with $e=0$ ($\mathrm{III}$: $5{,}534$ curves, $\mathrm{I}_0^*$: $5{,}394$) have minimum $0$, and for $\mathrm{I}_0^*$ the term runs up to $e=6$ (minimum $6$). Multiplicative primes show no Euler term: for curves with a single multiplicative prime $q$ with $\ell\mid n$ and no additive prime, $\min(v_\ell(\degphi)-v_\ell(n))=0$ for every value of $v_\ell(q^2-1)$ up to $6$. For two primes of potentially good type the terms add: for the $2{,}502$ working-set curves with exactly two odd primes of type $\mathrm{II}$ or $\mathrm{II}^*$ and no other contribution, $9\mid\degphi$ always, and the mechanism is not a congruence between $f$ and its own quadratic twist by the character ramified at both primes, which never occurs ($0$ of $2{,}502$).

\begin{table}[ht]\centering\footnotesize
\begin{tabular}{l rr rr rr rr}
\toprule
 & \multicolumn{2}{c}{$e=1$} & \multicolumn{2}{c}{$e=2$} & \multicolumn{2}{c}{$e=3$} & \multicolumn{2}{c}{$e=4$}\\
type ($\ell=3$) & min & $n$ & min & $n$ & min & $n$ & min & $n$\\
\midrule
$\mathrm{II}$ & 1 & 8,327 & 2 & 821 & 3 & 59 & 4 & 2\\
$\mathrm{II}^*$ & 1 & 4,184 & 2 & 296 & 3 & 12 & 4 & 1\\
$\mathrm{IV}$ & 1 & 4,290 & 2 & 296 & 3 & 12 & 4 & 1\\
$\mathrm{IV}^*$ & 1 & 8,221 & 2 & 821 & 3 & 59 & 4 & 2\\
$\mathrm{III}$ & 1 & 1,271 & 2 & 76 & 3 & 9 & -- & --\\
$\mathrm{III}^*$ & 1 & 1,271 & 2 & 76 & 3 & 9 & -- & --\\
$\mathrm{I}_0^*$ & 1 & 1,370 & 2 & 4,432 & 3 & 2,737 & 4 & 247\\
$\mathrm{I}_n^*$ & 1 & 30,497 & 2 & 2,589 & 3 & 145 & 4 & 13\\
\bottomrule
\end{tabular}

\caption{Working set, $\ell=3$: minimum of $v_3(\degphi)$ (number of curves) over curves with $E[3]$ irreducible whose only additive prime is a single $q\ge5$ of the given type with $e_3(E,q)=e$, and no multiplicative $3$-contribution. The same holds at $\ell=5$ and $7$.}\label{tab:iso}
\end{table}

\subsection{The two corollary forms}\label{sec:forms}
Table~\ref{tab:counts} gives, for each range and each small prime $\ell$, the number of optimal non-CM curves with $E[\ell]$ irreducible and a positive right-hand side in \eqref{eq:main} and in \eqref{eq:twist}, the number of violations of each inequality, and the number of Eisenstein curves with a positive bound together with the violations of Conjecture~\ref{conj:eis}. There are none. The theorem range of Theorem~\ref{thm:known} covers $281{,}739$ of the $913{,}244$ working-set pairs $(E,\ell)$ with $E[\ell]$ irreducible and a positive multiplicative bound; the remaining pairs have $\ell=3$ ($455{,}707$), $\ell\mid N$ ($175{,}798$ at $\ell\ge5$), or a contributing prime $q\equiv\pm1\pmod\ell$ ($410{,}074$), and show no violation either. Among working-set curves with a positive bound, equality holds in \eqref{eq:twist} for $39.4\%$ at $\ell=3$, $71.1\%$ at $\ell=5$, $81.2\%$ at $\ell=7$ and $91.5\%$ at $\ell=13$, and in \eqref{eq:main} for $36.9\%$, $69.7\%$, $80.4\%$ and $91.2\%$; the rule raises these to the rates of Section~\ref{sec:sharp} and bounds $861{,}074$ curves instead of $705{,}563$ at $\ell=3$.

\begin{table}[ht]
\centering\footnotesize
\begin{tabular}{llrrrrrr}
\hline
set & $\ell$ & curves & \eqref{eq:main} $>0$ & \eqref{eq:twist} $>0$ & viol.\ \eqref{eq:main} & viol.\ \eqref{eq:twist} & Eis.\ $>0$ / viol. \\
\hline
working set, $N \le 300000$ & 3 & 1,311,066 & 598,833 & 705,563 & 0 & 0 & 67,845 / 0 \\
 & 5 & 1,311,066 & 235,510 & 262,531 & 0 & 0 & 4,454 / 0 \\
 & 7 & 1,311,066 & 121,947 & 136,707 & 0 & 0 & 759 / 0 \\
 & 11 & 1,311,066 & 42,334 & 47,040 & 0 & 0 & 0 / 0 \\
 & 13 & 1,311,066 & 27,215 & 29,943 & 0 & 0 & 26 / 0 \\
\hline
hold-out, $300000 < N \le 400000$ & 3 & 429,936 & 206,749 & 242,691 & 0 & 0 & 18,225 / 0 \\
 & 5 & 429,936 & 79,550 & 89,164 & 0 & 0 & 989 / 0 \\
 & 7 & 429,936 & 41,387 & 46,767 & 0 & 0 & 137 / 0 \\
 & 11 & 429,936 & 14,244 & 15,970 & 0 & 0 & 0 / 0 \\
 & 13 & 429,936 & 9,212 & 10,280 & 0 & 0 & 2 / 0 \\
\hline
extension, $400000 < N \le 500000$ & 3 & 359,009 & 181,075 & 214,033 & 0 & 0 & 9,894 / 0 \\
 & 5 & 359,009 & 70,514 & 78,650 & 0 & 0 & 474 / 0 \\
 & 7 & 359,009 & 36,734 & 41,317 & 0 & 0 & 70 / 0 \\
 & 11 & 359,009 & 13,347 & 14,975 & 0 & 0 & 0 / 0 \\
 & 13 & 359,009 & 8,467 & 9,465 & 0 & 0 & 5 / 0 \\
\hline
out of table, $500000 < N \le 4\cdot 10^6$ & 3 & 893 & 561 & 718 & 0 & 0 & 0 / 0 \\
 & 5 & 893 & 6 & 7 & 0 & 0 & 0 / 0 \\
 & 7 & 893 & 0 & 1 & 0 & 0 & 0 / 0 \\
\hline
\end{tabular}

\caption{Verification counts for the two corollary forms. ``\eqref{eq:main} $>0$'' is the number of curves with $E[\ell]$ irreducible, not CM, for which the right-hand side of \eqref{eq:main} is positive, and ``viol.'' the number of violations; likewise for \eqref{eq:twist}. The last column counts the Eisenstein curves with a positive bound in \eqref{eq:main} and the violations of Conjecture~\ref{conj:eis}.}\label{tab:counts}
\end{table}

\subsection{The type \texorpdfstring{$\mathrm{III}$}{III} anomaly}\label{sec:anomaly}
With the $\mathrm{III}$ terms included in \eqref{eq:rule}, there are $364$, $68$, $43$ and $0$ exceptions on the four ranges, all at $\ell=3$, all with shortfall exactly one, and all of one shape: a prime $q$ of type $\mathrm{III}$ or $\mathrm{III}^*$ whose Euler factor is nontrivial with Frobenius acting trivially on the invariant line ($q\equiv1\pmod{12}$ for $q\ge5$; at $q=2$ the factor $1+2X$, which occurs with $v_2(N)=8$), together with a prime $p\equiv-1\pmod3$ with $t_3(E,p)\ge1$ (a multiplicative prime, or one of type $\mathrm{I}_n^*$ with $3\mid n$). On the clean configurations of the working set (one such $\mathrm{III}$ prime at $q\ge5$, one such $p$, nothing else; $1{,}066$ curves) the failure rate is $2\%$; it is $0$ for $q\equiv11\pmod{12}$ ($646$ curves), and it rises with $v_3(q-1)$ and with $v_3(p+1)$, to $12\%$ for $p\equiv8\pmod9$; it does not depend on whether the multiplicative reduction at $p$ is split. For every other type the Euler term adds to the Tamagawa terms with no exception, in particular when all the Tamagawa primes are $\equiv-1\pmod3$ ($14{,}000$ curves). Type $\mathrm{III}$ is the only type whose residual inertia image mod $3$ is semisimple and non-scalar; in the language of Conjecture~\ref{conj:selmer}, the exceptions are the cases where the dual Selmer correction is $1$. The configuration, a Steinberg prime $p\equiv-1\pmod\ell$, is the one excluded by the hypotheses of Kim and Ota, and the one in which a Taylor--Wiles level-raising argument sees an extra local term.

\subsection{The wild primes and the prime \texorpdfstring{$\ell$}{l}}\label{sec:wild}
Table~\ref{tab:isolated}, kept from the earlier version, isolates the additive types at $\ell=3$ with the weights of \eqref{eq:twist}: for the untwisted types it lists, for working-set curves with no multiplicative prime $q$ with $3\mid n_q$ and all additive primes of a single Kodaira type, the minimum of $v_3(\degphi)$ by type and number of primes; for the twisted types the minimum over the curves whose only contribution to \eqref{eq:twist} comes from primes of the given type at odd $q$, or at $q=2$, with the given total weight. The rows show that the prime $2$ contributes nothing through the types $\mathrm{II}$, $\mathrm{II}^*$ ($v_3(\degphi)=0$ for $14{,}182$ of the $44{,}672$ curves whose only candidate is such a prime) and nothing through a type $\mathrm{I}_n^*$ with $3\,\|\,n$, although with $9\mid n$ the minimum is $1$. Definition~\ref{def:terms} with PARI's factor explains every one of these minima: a type $\mathrm{II}$ at $2$ has $P_2=1$ and $e_3=0$; a type $\mathrm{IV}$ or $\mathrm{IV}^*$ at $2$ has $e_3=1$; a type $\mathrm{I}_n^*$ at $2$ has $t_3=v_3(-v_2(j))$ when it is potentially multiplicative and $0$ otherwise, and the index $n$ of the symbol is not the valuation of the Tate parameter at $2$, which is why $n$ was not the right invariant. With the weights of \eqref{eq:twist} at $2$ counted as if $2$ were odd, \eqref{eq:twist} fails thousands of times in each range; with the terms of Definition~\ref{def:terms} it never does. At $q=3$ with potentially good reduction and $\ell=3$, the minima of $v_3(\degphi)$ over curves with a single wild additive prime are $1$ for $\mathrm{II},\mathrm{IV},\mathrm{III}^*$, $2$ for $\mathrm{II}^*,\mathrm{IV}^*$ and $0$ for $\mathrm{III},\mathrm{I}_0^*$: these are the weights $w_3$ of Conjecture~\ref{conj:euler}, which we cannot derive from an Euler factor. For $\ell=5,7$ no term is needed at $q=\ell$ or at $q=2,3$ beyond Definition~\ref{def:terms}.

\begin{table}[ht]
\centering\small
\begin{tabular}{lccc}
\hline
type & weight 1 & weight 2 & weight 3 \\
\hline
IV (any $q$) & 1 (25,414) & 2 (969) & 3 (11) \\
IV$^*$ (any $q$) & 1 (30,168) & 2 (2,268) & 3 (45) \\
III (any $q$) & 0 (34,904) & 0 (2,208) & 0 (34) \\
III$^*$ (any $q$) & 0 (25,574) & 0 (1,307) & 0 (29) \\
$\mathrm{I}_0^*$ (any $q$) & 0 (42,924) & 0 (3,093) & 0 (53) \\
II, II$^*$ at odd $q$ & 1 (56,509) & 2 (2,502) & 3 (12) \\
$\mathrm{I}_n^*$ at odd $q$, weight $v_3(n)$ & 1 (28,175) & 2 (3,453) & 3 (89) \\
II, II$^*$ at $q=2$ & 0 (44,672) & -- & -- \\
$\mathrm{I}_n^*$ at $q=2$, weight $v_3(n)$ & 0 (21,176) & 1 (5,329) & 1 (342) \\
\hline
\end{tabular}

\caption{Minimum of $v_3(\degphi)$ (number of curves) over working-set curves with $E[3]$ irreducible whose only possible contributions at $3$ come from the primes named in the row, with the given number of primes (untwisted rows) or total weight of \eqref{eq:twist} (twisted rows).}\label{tab:isolated}
\end{table}

\subsection{Out of table}\label{sec:oot}
To test the statements beyond the tables, we generated curves with prescribed local types and conductor between $500000$ and $4\cdot10^6$ (families with a type $\mathrm{IV}$ or $\mathrm{IV}^*$ prime, with two or three type $\mathrm{II}$ primes, with mixed types, and random curves), kept those with trivial isogeny class (so that the curve is optimal and $E[\ell]$ is irreducible for every $\ell$) and not CM, and computed their modular degrees with \texttt{ellmoddegree}. Table~\ref{tab:oot} summarises them; Conjecture~\ref{conj:euler} and its two corollary forms hold for all of them (Tables~\ref{tab:uniform} and~\ref{tab:counts}, last rows), with equality in \eqref{eq:rule} for $51.6\%$ of the $827$ curves with a positive bound at $\ell=3$.

\begin{table}[ht]
\centering\small
\begin{tabular}{lrrr}
\hline
family & curves & largest conductor & median seconds per degree \\
\hline
II,II & 98 & 3,968,800 & 0.4 \\
II,II,II & 1 & 2,832,200 & 0.2 \\
IV & 334 & 3,988,012 & 1.1 \\
IV* & 104 & 3,994,600 & 4.9 \\
IVmixed & 266 & 3,983,504 & 1.3 \\
random & 90 & 2,710,716 & 32.1 \\
all & 893 & 3,994,600 & 1.5 \\
\hline
\end{tabular}

\caption{Curves outside the tables, with modular degrees computed independently.}\label{tab:oot}
\end{table}

\subsection{The Eisenstein case}\label{sec:eis}
Among the $75{,}430$ working-set optimal non-CM curves with a rational $3$-isogeny, $67{,}845$ have a positive bound in \eqref{eq:main} and $2{,}777$ a positive deficit $D$; $D=2$ occurs $22$ times. No curve without a rational $3$-torsion point has a deficit, whatever its $3$-isogeny structure. With a rational point of order $3$ and mod $3$ image of Borel type, $D\le1$; $D=2$ occurs only when $E[3]\simeq\Z/3\oplus\mu_3$ (the split Cartan image \texttt{3Cs.1.1} in the notation of \cite{Sutherland16}, seventeen curves, all of rank $0$ and torsion $\Z/3$) or when $E$ has a point of order $9$ (five curves). The two sources of a deficit of $2$ never combine: none of the $2{,}065$ curves of conductor at most $400000$ with mod $3$ image \texttt{3Cs.1.1} has a point of order $9$ (their torsion is $\Z/3$ or $\Z/6$), and the $18$ optimal curves with a point of order $9$ all have image \texttt{3B.1.1}. We have not proved that the two cannot coexist, so the bound $D\le2$ is part of the conjecture. The bounds of Conjecture~\ref{conj:eis} hold with no exception on all three ranges, and the torsion bound is attained with $D>0$ in $2{,}773$ cases. For $\ell=5$ and $7$ the deficits are at most $1$ and occur only with a rational $\ell$-torsion point ($192$ and $18$ cases in the working set). Measured against the stronger bound \eqref{eq:twist} the deficit is not controlled by the torsion: it exceeds the bound of Conjecture~\ref{conj:eis} for $1{,}148$ working-set curves at $\ell=3$ ($189$ and $39$ on the hold-out and the extension) and for $141$ at $\ell=5$. In the Eisenstein case the additive terms do not contribute reliably, and Conjecture~\ref{conj:eis} is stated for \eqref{eq:main} only.

\section{The prime 2}\label{sec:two}
At $\ell=2$ the modular degree is the subject of Watkins's conjecture $2^{\operatorname{rank}E(\Q)}\mid\degphi$ \cite{Watkins02} and of a substantial literature \cite{Dummigan06,CalegariEmerton09,Yazdani11,CaroPasten22,HatleyKundu,BKP}, and the mod $2$ representation has a rational $2$-isogeny for a large fraction of curves. We record only the following. For working-set optimal curves without a rational $2$-isogeny and not CM, the inequality $v_2(\degphi)\ge\sum_{q\,\|\,N}v_2(n_q)+\#\{q:\ \text{type }\mathrm{III},\mathrm{III}^*,\mathrm{I}_0^*\text{ or }\mathrm{I}_n^*\}$ holds with no exception on all three ranges ($964{,}432$, $325{,}510$ and $313{,}741$ curves), while counting $\mathrm{I}_0^*$ with weight $2=v_2(\#\Phig_q)$ fails ($123$, $26$ and $20$ exceptions). The isolated-configuration minima suggest that types $\mathrm{III}$, $\mathrm{III}^*$ and $\mathrm{I}_n^*$ contribute $3$ each and types $\mathrm{II}$, $\mathrm{IV}$ and $\mathrm{I}_0^*$ contribute $1$ each, in line with the theorem of Calegari and Emerton that a curve of odd modular degree has at most two odd bad primes and a rational $2$-torsion point or prime conductor \cite{CalegariEmerton09}. We have not attempted to reconcile the exponents with the results of \cite{Dummigan06,CaroPasten22,HatleyKundu,BKP}, nor with the Euler factors at $\ell=2$.

\section{Open questions}
\begin{enumerate}
\item Prove Conjecture~\ref{conj:selmer} for the types $\mathrm{II}$, $\mathrm{IV}$ and their duals, that is, that the $\Sigma$-imprimitive adjoint Selmer group of \cite{DFG04} has length at least the local terms; and identify the dual Selmer correction that makes the type $\mathrm{III}$ term fail by one in the configuration of Section~\ref{sec:anomaly}, which is the configuration excluded by the hypotheses of \cite{KimOta}.
\item Remove the restrictions $\ell\ge5$ and $q\not\equiv\pm1\pmod\ell$ from the Tamagawa terms (Theorem~\ref{thm:C}) and $\ell\nmid N$ from all the theorems; the data say the inequality survives all of them, and whether the equality \eqref{eq:ko} of Kim and Ota does is a separate question. The proof of Theorem~\ref{thm:euler} allows $\ell=3$, so the Tamagawa terms are the only place where $\ell=3$ is still open for the types $\mathrm{I}_n$ and $\mathrm{I}_n^*$.
\item Identify the weights $w_3$ at $q=\ell=3$ (one for $\mathrm{II},\mathrm{IV},\mathrm{III}^*$, two for $\mathrm{II}^*,\mathrm{IV}^*$, nothing for $\mathrm{III},\mathrm{I}_0^*$) with the local term at $\ell$ of the congruence-ideal formula, where the Euler factor is replaced by a condition at $\ell$; and explain the wild factors at $2$ and $3$ conceptually, rather than through PARI.
\item Identify the deficit of Conjecture~\ref{conj:eis} with a natural invariant of the Eisenstein ideal at $\ell$, presumably through the cuspidal subgroup of $J_0(N)$ \cite{Mazur77}.
\item Reconcile Section~\ref{sec:two} with the literature on Watkins's conjecture and with the Euler factors at $\ell=2$.
\end{enumerate}

\appendix
\section{Reproducibility}
All scripts, the pinned data commit and the result tables are in the repository \url{https://github.com/dsr001001/mcptrial} (directory \texttt{ecmine}). The symmetric-square Euler factors are computed by \texttt{scripts/step10\_sym2\_euler.py}, the rule is checked on the whole ranges by \texttt{scripts/step10\_rule.py}, the two corollary forms by \texttt{scripts/conj\_check.py}; each runs on the full tables in minutes to hours. A single curve can be checked against \eqref{eq:rule} in PARI/GP (version 2.15) with the following function, which returns $[v_\ell(\degphi),\text{bound}]$; it is the file \texttt{scripts/euler\_rule.gp}.
{\footnotesize\begin{verbatim}
rule(E, l) =
{ my(R = ellglobalred(E), P = R[4][,1], K = R[5], S = lfunsympow(E, 2), b = 0, j = E.j);
  for(i = 1, #P,
    my(q = P[i], k = K[i][2], e = 0, t = 0, vj = valuation(j, q));
    if(q != l,
      my(F = 1/lfuneuler(S, q), r = subst(F, x, 1/q^2)); \\ P_q(q^{-2}), the omitted factor
      if(k >= 5, r /= (1 - 1/q^2));                     \\ multiplicative: remove the naive one
      if(abs(k) != 3, e = valuation(r, l)),             \\ III, III* (codes 3, -3) not counted
      if(l == 3, e = if(k == 4 || k == 2, 1,
                     if(k == -4 || k == -2, 2, if(k == -3, 1, 0)))));   \\ w_3 at q = 3
    if(vj < 0, t = valuation(-vj, l));                   \\ Tamagawa exponent of the adjoint
    b += e + t);
  [valuation(ellmoddegree(E), l), b]; }
\end{verbatim}}
Here \texttt{K[i][2]} is PARI's Kodaira code ($n+4$ for $\mathrm{I}_n$, $-(n+4)$ for $\mathrm{I}_n^*$, $\pm2$ for $\mathrm{II}/\mathrm{II}^*$, $\pm3$ for $\mathrm{III}/\mathrm{III}^*$, $\pm4$ for $\mathrm{IV}/\mathrm{IV}^*$, $-1$ for $\mathrm{I}_0^*$). For example \texttt{rule(ellinit([1,1,1,-30,-76]), 3)} (the curve 121a1) returns \texttt{[1, 1]}, and \texttt{rule(ellinit([1,0,1,170,-3237]), 3)} (539c1) returns \texttt{[2, 2]}; \texttt{ellmoddegree} needs \texttt{parisizemax} of a few hundred megabytes for conductors in the thousands. The pipeline was rerun from a clean clone and reproduced every count of Table~\ref{tab:counts}. Lemma~\ref{lem:twist}(ii) was also checked numerically, with the script \texttt{twist\_lemma\_check.gp} of the repository, on $4{,}577$ random pairs (curve, odd bad prime), with no exception, and Lemma~\ref{lem:euler} against \texttt{lfunsympow} on every curve of the tables.

\section*{Acknowledgements}
[[The computations, the literature search and a first draft of this note were carried out with the assistance of an AI system (Claude, Anthropic); the author is responsible for all statements and proofs.]]

\bibliographystyle{amsplain}
\bibliography{refs}
\end{document}
